\documentclass[11pt,a4paper]{article}
\usepackage[T1]{fontenc}
\usepackage[utf8]{inputenc}
\usepackage{lmodern,amsmath,amssymb,amsthm,mathtools,booktabs}
\usepackage[margin=25mm]{geometry}
\usepackage{microtype}
\usepackage[colorlinks=true,linkcolor=blue,citecolor=blue,urlcolor=blue]{hyperref}
\hypersetup{pdftitle={Adaptive estimation of rare signs with a known finite gap law},
 pdfsubject={A closed quantile characterization under expected observation budgets}}
\newtheorem{theorem}{Theorem}[section]
\newtheorem{lemma}[theorem]{Lemma}
\newtheorem{proposition}[theorem]{Proposition}
\newtheorem{corollary}[theorem]{Corollary}
\theoremstyle{remark}
\newtheorem{remark}[theorem]{Remark}
\newcommand{\E}{\mathbb E}
\newcommand{\Pp}{\mathbb P}
\newcommand{\R}{\mathbb R}
\newcommand{\1}{\mathbf 1}
\newcommand{\Q}{\mathcal Q}
\DeclareMathOperator{\Var}{Var}
\DeclareMathOperator{\KL}{KL}
\DeclareMathOperator{\TV}{TV}
\DeclareMathOperator{\Ent}{Ent}
\DeclareMathOperator{\Aff}{Aff}
\DeclareMathOperator{\kl}{kl}
\setlength{\parindent}{0pt}
\setlength{\parskip}{.35em}
\setlength{\emergencystretch}{2em}
\setcounter{tocdepth}{2}
\title{\bfseries Adaptive estimation of rare signs\\with a known finite gap law}
\author{}
\date{29 October 2026}
\begin{document}
\maketitle
\begin{abstract}
Each fresh coin has an unobserved, persistent Bernoulli bias
$1/2+SD$, where the magnitude $D$ has a known finite law and the
conditional positive-sign probability is unknown. We characterize the
minimax root-mean-square error for estimating the population fraction
of positive signs, bounded by a known rarity parameter $\rho$, under a
uniform expected observation budget. Algorithms may create, revisit
and interleave coins adaptively. The rate has a closed form: remove
the smallest-gap mass $\rho$, retain its truncated inverse-gap cost,
and add the square root of the dyadic inverse-square cost entropy
of the remaining mass, scaled by rarity and budget. The lower bound
combines a posterior-martingale tail obstruction, entropy packing,
a one-coin activation argument and a common-pair reduction valid for
fully adaptive population algorithms. The upper bound uses an
observable likelihood-ratio classifier race and a finite multilevel
Poisson allocation. All specialized ingredients are proved, and
explicit universal comparison constants are given.
\end{abstract}
\section{Statistical model and main theorem}\label{sec:model}

\subsection{Persistent coins and the observation budget}
Let
\[
 \mu=\sum_{i=1}^n w_i\delta_{d_i},\qquad
 0<d_1<\cdots<d_n\le\tfrac12,\qquad
 w_i>0,\quad\sum_{i=1}^n w_i=1,
\]
be a known finite probability law. Equal gaps are merged before using
this representation. A parameter is a function $q$ on the support of
$\mu$. The complete latent pairs $(D_k,S_k)$ of distinct fresh coins
are iid, with $D_k\sim\mu$, $S_k\in\{-1,+1\}$, and
$\Pp_q(S_k=+1\mid D_k=d)=q(d)$. Each coin retains its pair on every
revisit. Conditional on all latent pairs, the observation tapes of
different coins are independent, and tape $k$ consists of iid
Bernoulli$(1/2+S_kD_k)$ observations. Neither $D_k$ nor $S_k$ is observed.

Fix a known rarity bound $0<\rho\le1/4$. The parameter class and target are
\begin{equation}\label{eq:parameter-class}
 \mathcal Q_\rho=\left\{q:\ 0\le q(d_i)\le1,\quad
       \theta(q):=\sum_{i=1}^n w_iq(d_i)\le\rho\right\}.
\end{equation}
This is a global constraint on the positive-sign fraction; it does not
impose $q(d_i)\le\rho$ at each gap.

An algorithm may create, revisit and interleave coins, and use an
independent random seed whose distribution does not depend on $q$.
Its choices and output are measurable functions of that seed, coin
identities and purchased observations. It must terminate and return
a real estimate $\widehat\theta$ almost surely under every
$q\in\mathcal Q_\rho$. Each observed flip costs one, including the
first observation of a fresh coin. Let $N$ be the total number of
observations. Coin creation, arithmetic and random-number generation
are not charged; creating an unobserved coin yields no information.
Zero observations are allowed, and no deterministic total cap is imposed.

For a known budget $B>0$, define the minimax root-mean-square error
\begin{equation}\label{eq:minimax-risk}
 R_{\mu,\rho}(B)=
 \inf_{\mathcal A:\ \sup_{q\in\mathcal Q_\rho}\E_qN\le B}
 \ \sup_{q\in\mathcal Q_\rho}
       \left\{\E_q\bigl(\widehat\theta-\theta(q)\bigr)^2\right\}^{1/2}.
\end{equation}
The budget is uniform over the entire parameter class. A procedure may
depend on $(\mu,\rho,B)$, but not on $q$ or any genuine coin's latent
state. Projecting its estimate onto $[0,\rho]$ cannot increase loss
or work. We abbreviate $R_{\mu,\rho}(B)$ to $R(B)$ when the inputs
are fixed.

\subsection{The closed quantile functional}
Remove exactly the mass $\rho$ with smallest gaps. In terms of the
weights, set
\begin{equation}\label{eq:quantile-split}
 \tau_i=\min\left\{w_i,\left(\rho-\sum_{\ell<i}w_\ell\right)_+\right\},
 \qquad \nu=\sum_i\tau_i\delta_{d_i},\qquad \kappa=\mu-\nu.
\end{equation}
Here $x_+=\max\{x,0\}$. Let $t$ be the smallest support point for
which $\mu((0,t])\ge\rho$. Then $\nu$ has mass $\rho$ and is supported
on $(0,t]$, while $\kappa$ has mass $1-\rho$ and is supported on
$[t,1/2]$. The atom at $t$ may be split fractionally.

For $j=0,1,2,\ldots$, define the dyadic bands and exact core costs by
\begin{equation}\label{eq:core-entropy}
 \begin{gathered}
 b_j=2^{-j-1},\qquad I_j=(b_j/2,b_j],\qquad
 a_j^*=\sum_{i:d_i\in I_j}\frac{w_i-\tau_i}{d_i^2},
 \qquad A^*=\sum_j a_j^*,\\
 \Gamma_\kappa=A^*+\sum_{j:a_j^*>0}a_j^*\log(A^*/a_j^*).
 \end{gathered}
\end{equation}
All logarithms are natural. There are finitely many positive costs,
and $A^*>0$. Thus $\Gamma_\kappa=A^*(1+\Ent(a^*/A^*))$ is the
inverse-square cost of the core multiplied by one plus the entropy
of its distribution across dyadic scales.

With $h=B^{-1/2}$, set
\begin{equation}\label{eq:closed-functional}
 T_h=\sum_i\tau_i\min\{1,h/d_i\},\qquad
 \Q(B)=\min\left\{\rho,\ T_h+h\sqrt{\rho\Gamma_\kappa}\right\}.
\end{equation}
The expression depends only on the known weights, gaps, rarity and
budget. No optimization over subsets or submeasures is needed to
evaluate it.

\subsection{The theorem and its quantifiers}
\begin{theorem}[Closed quantile characterization]\label{thm:direct-main}
For every integer $n\ge1$, every finite law $\mu$ described above,
every $0<\rho\le1/4$, and every $B>0$,
\begin{equation}\label{eq:main-bounds}
       10^{-20}\Q(B)\ \le\ R_{\mu,\rho}(B)\ \le\ 30001\Q(B).
\end{equation}
More precisely, for each $(\mu,\rho,B)$ there is one observable
algorithm satisfying, simultaneously for all $q\in\mathcal Q_\rho$,
\[
 \E_qN\le B,\qquad
 \left\{\E_q(\widehat\theta-\theta(q))^2\right\}^{1/2}
 \le30001\Q(B).
\]
Conversely, every algorithm with
$\sup_{q\in\mathcal Q_\rho}\E_qN\le B$ has worst-case RMS error
at least $10^{-20}\Q(B)$, including algorithms with arbitrary
adaptive interleaving, nonlinear outputs and unbounded stopping.
\end{theorem}

The constants are universal and are not optimized. The claim is a
comparison up to constants, not an exact formula for the minimax
coefficient. The iid assumption concerns the complete latent pairs,
not merely their magnitudes. The theorem concerns known finite laws
and expected observation budgets; other information structures and
budget conventions are not included.

\subsection{Notation and testing conventions}\label{sec:notation}
Write $s=\mathbf1_{\{S=+1\}}\in\{0,1\}$ for a binary sign label;
thus the conditional observation mean is $1/2+(2s-1)d$. A
single-unit procedure means a procedure on one genuine coin, with
its latent state retained throughout. A null prior is $q\equiv0$;
$\E_0$ denotes expectation at that prior. When the two signs are
compared at one specified gap, $\E_{-d}$ and $\E_{+d}$ refer to
Bernoulli means $1/2-d$ and $1/2+d$, respectively. Fair-coin
expectations are explicitly labelled $\E_{\rm fair}$.

For probability laws $P,Q$, dominated by a common measure, write
\[
 \Aff(P,Q)=\int\sqrt{dP\,dQ},\qquad
 \mathfrak h(P,Q)=1-\Aff(P,Q),\qquad
 \TV(P,Q)=\sup_E|P(E)-Q(E)|.
\]
In particular, $\mathfrak h$ denotes squared Hellinger distance in
the convention used here, whereas $h=B^{-1/2}$ is a budget scale.
We use $\TV\le\sqrt{2\mathfrak h}$, $\mathfrak h\le\KL/2$ and
Pinsker's inequality $\TV\le\sqrt{\KL/2}$, as well as KL data
processing and the finite-prefix chain rule. Unbounded stopping
is handled by finite caps and the stated limiting arguments.
For a probability vector $p$, $\Ent(p)=\sum_{p_j>0}p_j\log(1/p_j)$;
$\log_+x=\max\{0,\log x\}$, with $\log_+0=0$. The symbols
$\lesssim$, $\gtrsim$ and $\asymp$ suppress universal positive
constants only.

\begin{lemma}[Two-point RMS testing]\label{lem:two-point}
If two admissible parameters have target separation $\Delta>0$
and terminal transcript laws of total variation at most $1/2$,
every estimator has maximum RMS error at those parameters at
least $\Delta/4$.
\end{lemma}
\begin{proof}
Threshold the estimate at the midpoint of the two target values.
Each testing error incurs squared estimation loss at least
$\Delta^2/4$. The sum of the two testing errors is at least
$1-\TV\ge1/2$. Hence the sum of the two mean-square errors is
at least $\Delta^2/8$, and their maximum is at least $\Delta^2/16$.
\end{proof}

\section{Proof strategy and entropy preliminaries}\label{sec:overview}
The two terms in \eqref{eq:closed-functional} describe different
obstructions. The removed mass $\nu$ may concentrate at gaps too
small to resolve at the current budget. Its contribution is the
truncated inverse-gap cost $T_h$. In the remaining mass $\kappa$,
collective sign uncertainty and identifying the relevant scales
produce the cost-entropy term $h\sqrt{\rho\Gamma_\kappa}$.

\subsection{Organization of the proof}
The lower proof in Section~\ref{sec:matching-lower} first obtains
$R(B)\ge T_h/32$ by a posterior-martingale calculation and an
adaptive midpoint decomposition. It then filters the dyadic core
costs. Either the tail already controls the desired rate, or the
filter preserves a constant fraction of the core entropy. Grouping
and trimming the surviving bands produces conditional experiments
that satisfy the required rarity and saturation constraints.
Their lower bounds add through a common-pair argument.

The upper proof in Section~\ref{sec:direct-upper} uses the same
quantile split as a known design choice. A likelihood bank races
an ordinary capped sign classifier on observations of the same
coin. Coupled corrections at successive accuracies are replicated
with independent Poisson counts. The resulting finite multilevel
estimator has uniform expected work at most $B$. A homogeneous
kernel estimate controls mixed work across scales.
Section~\ref{sec:assembly-contract} assembles all branches and
tracks the constants in Theorem~\ref{thm:direct-main}.

Appendix~\ref{sec:uf-foundations} proves the classification, work
and allocation bounds. Appendix~\ref{sec:lower-foundations}
proves weighted sign testing, the activation gate and the annular
one-coin obstruction. Appendix~\ref{sec:transfer-foundations}
proves finite convex comparison, the common pair for all capped
local policies, the adaptive midpoint decomposition and additive
localization. No statistical selection theorem is assumed.

\subsection{Relation to prior methods}
Lee and Valiant~\cite{lv} study adaptive estimation for mixtures
of unknown coins and provide a midpoint decomposition for
fully adaptive observations. We reproduce the finite-prefix
decomposition needed here, including observable types and the
passage to uniform expected budgets. Their sample-complexity
theorem is not used as a substitute for the RMS statement above.

The weighted sign-testing argument follows the gap-entropy
mechanism of Chen, Li and Qiao~\cite{clq}, including their
Bernoulli extension. The specialized separated-gap lemma is
proved in Appendix~\ref{sec:lower-foundations}. The finite
convex comparison has methodological antecedents in
Juditsky and Nemirovski~\cite{jn} and Polyanskiy and Wu~\cite{pw};
its exact version and the simultaneous common-pair construction
are proved in Appendix~\ref{sec:transfer-foundations}.
We use Hoeffding's inequality~\cite{hoeffding} and Sion's convex
minimax theorem~\cite{sion}, together with the standard probability
and compactness facts indicated where applied. The resulting
quantile formula and all specialized reductions are established
within this manuscript. These attributions identify methodological
dependencies and do not assert a comprehensive novelty comparison.

\subsection{Cost entropy and canonical band costs}\label{sec:entropy-preliminaries}
For a finite nonnegative vector $a$, put $A=\sum_j a_j$ and define
$\mathcal G(a)=A+\sum_{j:a_j>0}a_j\log(A/a_j)$, with
$\mathcal G(0)=0$. This continuous function is homogeneous,
coordinatewise increasing and concave. On the positive coordinates,
its derivative is $1+\log(A/a_j)$ and its Hessian is
$A^{-1}\mathbf1\mathbf1^\top-\operatorname{diag}(1/a_j)$.
Cauchy--Schwarz makes the latter negative semidefinite, and continuity
extends the inequalities to zero coordinates. Homogeneity and concavity
imply superadditivity, and binary entropy gives
\begin{equation}\label{eq:binary-mixture}
 \mathcal G(a+b)\le(1+\log2)[\mathcal G(a)+\mathcal G(b)]
                 \le2[\mathcal G(a)+\mathcal G(b)].
\end{equation}
Indeed, keep and then forget the label specifying which cost vector
generated a coordinate. The extra entropy is at most $\log2$ times
the total cost, which is at most $\mathcal G(a)+\mathcal G(b)$.

In the lower bound we use canonical costs
\[
 v_j=\kappa(I_j),\quad a_j=v_j/b_j^2,\quad A=\sum_j a_j,
 \quad F=\mathcal G(a),\quad \lambda=\sqrt{\rho/F}.
\]
The core has mass $1-\rho>0$, so $F>0$. Coordinatewise comparison yields
\begin{equation}\label{eq:canonical-core}
                      F\le\Gamma_\kappa\le4F.
\end{equation}
Only costs are rounded; the observation experiment retains its actual gaps.


\section{Lower bound}\label{sec:matching-lower}
Throughout the proof, $h=B^{-1/2}$ is the budget scale and
$\mathfrak h(P,Q)$ is the squared Hellinger distance defined in
Section~\ref{sec:notation}.

\subsection{An explicit posterior-martingale bound for the quantile tail}
Fix a gap $d$ and a capped one-coin policy. Under the equal sign prior,
let $U_k=\E[S\mid\text{seed},X_1,\ldots,X_k]$, with $X_i\in\{-1,1\}$
and $\E[X_i\mid S]=2dS$. Then $U_0=0$ and $U_k$ is a martingale.
For a current posterior mean $u$,
\[
 U_{k+1}=\frac{u+2dX_{k+1}}{1+2duX_{k+1}},\qquad
 \E[(U_{k+1}-u)^2\mid U_k=u]
 =\frac{4d^2(1-u^2)^2}{1-4d^2u^2}\le4d^2.
\]
At $d=1/2$, the posterior becomes absorbing at $\pm1$ after the first
observation; the undefined zero-over-zero expression there is interpreted
through the zero increment. For the policy's capped stopping time $N$,
\begin{equation}\label{eq:posterior-energy}
                         \E_{1/2}U_N^2\le4d^2\E_{1/2}N.
\end{equation}
The seed may be arbitrary and independent of the latent state.

Change the conditional positive-sign probability to
$1/2-\epsilon$ or $1/2+\epsilon$, $0\le\epsilon\le1/4$.
Relative to the midpoint transcript law the endpoint densities are
$1-2\epsilon U_N$ and $1+2\epsilon U_N$.
Since $1-\sqrt{1-x}\le x$ on $[0,1]$, (\ref{eq:posterior-energy}) gives
\begin{equation}\label{eq:explicit-tail-local}
 \mathfrak h(P_{1/2-\epsilon},P_{1/2+\epsilon})
 \le4\epsilon^2\E_{1/2}U_N^2
 \le16\epsilon^2d^2\E_{1/2}N.
\end{equation}

\begin{lemma}[Quantile-tail obstruction]\label{lem:tail}
For every allowed law, rarity and budget, $R_{\mu,\rho}(B)\ge T_h/32$.
\end{lemma}
\begin{proof}
Use an oracle experiment that lets the algorithm choose any tail atom
$i$ as an observable type, then obtain a genuine coin with that fixed
gap $d_i$ and unknown sign. Each flip is still charged. The algorithm
may interleave types and revisit coins.
Set
\[
 \epsilon_i=\frac1{16}\min\{1,h/d_i\},\qquad
 q_{i,\pm}=1/2\pm\epsilon_i.
\]
Equation (\ref{eq:explicit-tail-local}) bounds the local distance for
every type by $(h^2/16)$ times its midpoint work. The midpoint
decomposition, proved in Lemma~\ref{lem:tf-midpoint}, therefore bounds
global distance by $h^2B/16=1/16$. It permits observable types chosen
adaptively and arbitrary interleaving.

Every original algorithm can be simulated in this oracle experiment:
for each fresh coin privately choose tail type $i$ with probability
$\tau_i$; on the remaining mass simulate an all-negative coin with
the residual magnitude law $\kappa/(1-\rho)$ for free. Keep that
private choice on revisits. The resulting original endpoints are
\[
 q_\pm(d)=\frac{\mathrm d\nu}{\mathrm d\mu}(d)\left(\frac12\pm\frac1{16}\min\{1,h/d\}\right).
\]
The density uses total masses at a support point. At a distinct atom it
equals $\tau_i/w_i$; this remains well-defined even if an input lists
identical gaps more than once.
Both endpoints and their midpoint are admissible, since their total
positive mass is at most $(9/16)\rho$. The oracle's genuine work is
at most the original flip count, so its midpoint budget is at most $B$.
No budget assumption over arbitrary oracle parameters is needed.
The original target separation is
$2\sum_i\tau_i\epsilon_i=T_h/8$.
Data processing gives total variation at most
$\sqrt{2/16}<1/2$. Two-point midpoint testing gives RMS at least one
quarter of the separation, namely $T_h/32$.
For unbounded original stopping, cap the global experiment first.
The same endpoint pair works at every cap; increasing prefixes give
the terminal TV bound, as proved after Lemma~\ref{lem:tf-midpoint}.
\end{proof}

\begin{lemma}[Rare Bernoulli-label floor]\label{lem:floor}
$R_{\mu,\rho}(B)\ge\min\{\rho,h\sqrt\rho\}/32$.
\end{lemma}
\begin{proof}
Restrict to constant priors $q(d)=p$ and strengthen the experiment by
revealing a fresh coin's sign for unit cost. Its magnitude and subsequent
flips can then be simulated using independent randomness, so the
number of purchased labels is at most the original number of flips.
Take $p_- =\rho/2$, $p_+=p_-+\delta$ with
$\delta=\min\{\rho,h\sqrt\rho\}/8$.
One Bernoulli-label divergence is at most
$\delta^2/[p_+(1-p_+)]\le3\delta^2/\rho$.
The stopped chain rule bounds total KL by $3\delta^2B/\rho\le3/64$.
Pinsker gives TV below $1/2$, and midpoint testing gives RMS
at least $\delta/4$. Finite prefixes justify unbounded stopping.
\end{proof}

\subsection{A deterministic entropy filter}
Retain the canonical core notation $a,A,F,\lambda$ from
Section~\ref{sec:entropy-preliminaries}.
If $h\ge t$, then $T_h=\rho$, so the tail lemma already proves the
desired lower bound. Assume $h<t$. If $\lambda\ge t/4$, then
\[
 h\sqrt{\rho\Gamma_\kappa}\le2h\sqrt{\rho F}
 =2h\rho/\lambda\le8h\rho/t\le8T_h,
\]
and $\Q\le9T_h$. Thus it remains to treat $\lambda<t/4$.

Set $L_j=1+\log(A/a_j)$ on positive-cost bands. Declare a band
\emph{good} exactly when
\begin{equation}\label{eq:good-band}
                       b_j\ge\tfrac\lambda2 L_j.
\end{equation}
This is a direct test of the known weights and endpoints.
Let $F_g$ and $F_b$ be the entropies of the good and bad cost vectors,
each normalized by its own total cost.

\begin{lemma}[Filtering retains core entropy]\label{lem:filter}
If $\lambda<t/4$, then $F_b\le F/8$ and $F_g\ge F/4$.
\end{lemma}
\begin{proof}
Every core band has endpoint at least $t$. Number the bands upward
from the one containing $t$, with endpoints $\beta_0 2^k$, $k\ge0$.
On a bad band, $L_j>x_j:=2b_j/\lambda$, and $x_j\ge8$.
Since $a_j/A=e^{1-L_j}$ and $xe^{1-x}$ decreases for $x\ge1$,
\[
 \sum_{j\text{ bad}}a_jL_j
 \le A\sum_{k\ge0}x_0 2^k e^{1-x_0 2^k},\qquad x_0=2\beta_0/\lambda\ge8.
\]
The ratio of successive summands is at most $2e^{-8}<1/2$.
Their sum is at most $2x_0e^{1-x_0}\le16e^{-7}<1/8$.
The bad vector's own entropy is at most this full-normalization sum,
so $F_b\le A/8\le F/8$. By (\ref{eq:binary-mixture}),
$F\le2(F_g+F_b)$; consequently $F_g\ge3F/8\ge F/4$.
\end{proof}

\subsection{Packing the filtered entropy into actual experiments}
We first group coordinates with comparable costs.
\begin{lemma}[Entropy tiers]\label{lem:entropy-tiers}
For a finite positive vector $a$, let $M=\max_j a_j$ and group its
coordinates by $a_j\in(M2^{-k-1},M2^{-k}]$, $k\ge0$.
If $F_k$ is each group's entropy cost, then
\begin{equation}\label{eq:entropy-tiers}
 \mathcal G(a)\le7\sum_k F_k.
\end{equation}
\end{lemma}

\begin{proof}
Let $A_k$ be the cost mass in group $k$, $n_k$ its cardinality, and
$p_k=A_k/A$. Entropy decomposition gives
$\mathcal G(a)=\sum_kF_k+A\Ent(p)$.
Empty groups are omitted. For the comparison on $k\ge0$, the reference
probability is $q_k=2^{-k/2}/(2+\sqrt2)$, so $\KL(p\Vert q)\ge0$
applies even when the nonempty tier indices are sparse.
Comparison with the probability vector proportional to $2^{-k/2}$ gives
\[
 A\Ent(p)\le A\log(2+\sqrt2)
                 +\frac{\log2}{2}\sum_k kA_k.
\]
On groups with $n_k\ge2^{k/2}$, the second term is at most
$\sum_kA_k\log n_k\le\sum_kF_k$; comparable weights imply
$F_k\ge A_k[1+\log n_k-\log2]\ge A_k\log n_k$.
On the other groups, $A_k\le M2^{-k/2}$, so their contribution is at most
\[
 \frac{\log2}{2}M\sum_{k\ge0}k2^{-k/2}
 =\frac{(\log2)\sqrt2}{2(\sqrt2-1)^2}M<3M\le3A.
 \]
Finally $\log(2+\sqrt2)<2$ and $A\le\sum_kF_k$.
Adding the initial $\sum_kF_k$ proves \eqref{eq:entropy-tiers}.
\end{proof}


Split the good bands into high bands $j<6$ and low bands $j\ge6$,
with entropies $F_H,F_L$. Since $b_j\ge1/64$ on the first six bands,
$F_H\le4096(1+\log6)=:C_H$.
If $F_H\ge F_L$, then
\[
 F\le4F_g\le8(F_H+F_L)\le16C_H,
\]
and the Bernoulli-label floor implies
$R\gtrsim\min\{\rho,h\sqrt{\rho F}\}$.

Otherwise $F_L>F_H$, so $F\le16F_L$.
Split the low bands by index parity, then within each parity into the
comparable-cost tiers of Lemma~\ref{lem:entropy-tiers}.
The sum of the resulting group entropies is at least $F_L/14$.
Indeed $F_L\le2(F_{\rm even}+F_{\rm odd})$, and the tier lemma is
applied separately within each parity. Both parities and all nonempty
tiers are retained; choosing just one tier would not justify this bound.
Each group has endpoint ratios at least four and costs in $(c,2c]$
for some $c>0$.

For a group of $n$ bands, (\ref{eq:good-band}) and $A\ge nc$ give
\[
 b_j\ge\frac\lambda2[1+\log(n/2)]
       \ge\frac\lambda8(1+\log n).
\]
Remove its $s=\lceil\log_4n\rceil$ smallest-gap bands, with $s=0$
for $n=1$. For $n\ge2$, $s\le n/2$, so at least $n/2$ remain.
Let $G_i$ be the remaining group's entropy. Coordinate monotonicity yields
\begin{equation}\label{eq:trimmed-mass}
 G_i\ge F_{\rm group}/16,\qquad
 b_{\min,i}\ge\lambda n(1+\log n)/8.
\end{equation}
For the first bound, divide the lower bound
$cn'(1+\log n')$ by the upper bound $2cn(1+\log n)$;
the ratio is at least $(1-\log2)/4>1/16$.
The second bound follows because deleting $s$ smallest bands raises
the minimum endpoint by at least $4^s\ge n$.
The case $n=1$ is immediate. The integer inequality for $n\ge2$
can be checked directly at $n=2,3,4$; when
$4^{s-1}<n\le4^s$ and $s\ge2$, we have
$2s\le4^{s-1}+1\le n$. Trimming is by increasing gap (equivalently,
increasing endpoint), not by cost. The number $n$ in
\eqref{eq:trimmed-mass} is the pre-trim cardinality.

Let $\sigma_i$ be the restriction of the actual core measure to a
remaining group, $m_i$ its mass, $u_i$ its largest endpoint, and
$w_{*,i}$ its smallest band mass. Then
\[
 m_i/u_i^2\ge c,\qquad w_{*,i}\ge c b_{\min,i}^2,
 \qquad G_i\le2cn(1+\log n).
\]
Consequently
\begin{equation}\label{eq:direct-saturation}
                 \frac{m_iw_{*,i}}{u_i^2}\ge\lambda^2G_i^2/256.
\end{equation}
Assign original rarity $\theta_i=\lambda^2G_i/1024$ and conditional
rarity $r_i=\theta_i/m_i$.
Equation (\ref{eq:direct-saturation}) gives
$\theta_iG_i\le m_iw_{*,i}/(4u_i^2)$.
Because $G_i\ge m_i/u_i^2$, this actually implies
$\theta_i\le w_{*,i}/4\le m_i/4$, hence $r_i\le1/4$.
Here $w_{*,i}$ is a band mass, not a minimum individual atom mass.
Under the conditional law $\sigma_i/m_i$, the entropy and minimum
band mass are $G_i/m_i$ and $w_{*,i}/m_i$. Its entropy-branch condition is
\[
 r_i u_i^2(G_i/m_i)\le w_{*,i}/m_i
 \quad\Longleftrightarrow\quad
 \theta_iG_i\le m_iw_{*,i}/u_i^2.
\]
Thus the preceding factor-four margin verifies the actual conditional
condition. There is no unnormalized-to-conditional mass factor missing.
Disjointness and superadditivity give
\[
 \sum_i\theta_i\le\lambda^2F/1024=\rho/1024,
 \qquad \sum_iG_i\ge F_L/224\ge F/3584.
\]

All remaining bands have largest endpoint at most $1/128$ and
endpoint ratios at least four. Apply the annular one-unit obstruction
of Proposition~\ref{prop:scaled-annular} to each normalized component
$\sigma_i/m_i$, followed by the additive localization theorem
(Theorem~\ref{thm:additive-localization}). Their complete proofs are
included in Appendices~\ref{sec:lower-foundations} and
\ref{sec:transfer-foundations}.
The actual minimum-band-mass condition is verified by
(\ref{eq:direct-saturation}), so the entropy branch is active.
We obtain
\[
 R(B)\gtrsim\sum_i\min\{\theta_i,h\sqrt{\theta_iG_i}\}
 =\left(\sum_iG_i\right)
       \min\{\lambda^2/1024,h\lambda/32\}
 \gtrsim\min\{\rho,h\sqrt{\rho F}\}.
\]
The proportional allocation $\theta_i/G_i=\lambda^2/1024$ is what
makes this estimate valid jointly at every budget.

\subsection{Completion of the lower bound}
The cases $h\ge t$ and $\lambda\ge t/4$ were handled by the tail lemma.
In the remaining case, the floor or the packing argument gives
$R\gtrsim\min\{\rho,h\sqrt{\rho F}\}$, while always $R\ge T_h/32$.
Finally, by (\ref{eq:canonical-core}),
\[
 \Q(B)\le T_h+2\min\{\rho,h\sqrt{\rho F}\}.
\]
The two lower bounds control this sum up to a universal constant.
Thus $R(B)\ge c\Q(B)$, proving the lower half of
Theorem~\ref{thm:direct-main}. Section~\ref{sec:assembly-contract}
tracks the constants to obtain the stated value $10^{-20}$.

\section{Observable upper bound}\label{sec:direct-upper}
We construct an estimator at the quantile split from
Section~\ref{sec:model}. The classification and allocation lemmas
used below are proved in Appendix~\ref{sec:uf-foundations}.

\begin{theorem}[Upper bound at the fixed quantile split]\label{thm:direct-upper}
For every admissible $(\mu,\rho,B)$, there is an observable estimator
with uniform expected work at most $B$ and worst RMS error at most
$30001\Q(B)$.
\end{theorem}

If $\Q(B)=\rho$, use the constant estimator $\rho/2$.
Otherwise $h<t$, because $h\ge t$ would imply $T_h=\rho$.
Fix
\[
 \sigma=\kappa,\qquad \tau=\nu|_{[h,1/2]},\qquad
 \zeta=\nu|_{(0,h)},\qquad z=\zeta((0,1/2]),\quad
 H=\int d^{-1}\,\tau(dd).
\]
Then $\mu=\sigma+\tau+\zeta$ and $T_h=z+hH$.
These are analytical pieces only. The estimator never observes which
piece produced its genuine coin.
All likelihood evaluations and design parameters use the known law.
The complexity measure counts physical coin observations, including those
given to either side of the race; it does not assert a bound on arithmetic
time, memory, or random-number generation.

\subsection{An uncapped reference bank with a capped observable race}

Use the fixed decomposition above, put
$z=\zeta((0,1/2])$, and assume $\Gamma_\sigma,H_\tau<\infty$.
It suffices to treat $z\le\rho/128$: otherwise the constant $\rho/2$
has risk at most $64z$. For a finite submeasure $\xi$, define explicitly
\[
 a_j(\xi)=\int_{I_j}d^{-2}\,\xi(dd),\qquad
 A_\xi=\sum_j a_j(\xi),\qquad H_\xi=\int d^{-1}\,\xi(dd),
\]
\[
 \Gamma_\xi=A_\xi+\sum_j a_j(\xi)\log(A_\xi/a_j(\xi)).
\]
Zero coordinates are omitted, and all three quantities are zero for the
zero measure. These are exact inverse-square costs, not the rounded
canonical costs used in the lower bound. Write $A=A_\sigma$, $H=H_\tau$.
If $\sigma\ne0$, use all its occupied bands, with weights
$\pi_j=a_j(\sigma)/A$ and references $b_j/2$ and $b_j/4$.
For an internal allowance $u<\rho$, set $\epsilon=u/\rho$.
Stop negative if a negative reference likelihood against fair reaches
$1/(\epsilon\pi_j)$; stop positive if
\begin{equation}\label{eq:split-bank-ratio}
 \frac{\sum_j\pi_jP_{1/2+b_j/4,n}}
      {\int P_{1/2-d,n}\,\mu(dd)}\ge1/u.
\end{equation}
Give negative crossings priority. A zero denominator triggers a positive
stop. The known all-negative mixture in the denominator is the full law
$\mu$, regardless of the decomposition.

By Lemma~\ref{lem:uf-mixture-bank}, a wrong negative crossing has
conditional positive-sign probability at most $\epsilon$, and a positive
crossing has all-negative-mixture probability at most $u$. Thus the total
probability of a wrong bank stop is at most $2u$ over $\mathcal Q_\rho$.
Termination outside $\sigma$ is not assumed.

Race this bank against the ordinary classifier of Lemma~\ref{lem:uf-ordinary}, using
alternating observations of the same physical coin. At an ordinary cap
$M$, its conditional work certificate is
\[
 o_M(d)=\min\{M,17d^{-2}[1+\log_+(Md^2)]\}.
\]
For any finite measure $\xi$, put
\[
 \eta_\xi(M)=\int\min\{1,(Md^2)^{-2}\}\,\xi(dd),\qquad
 O_\xi(M)=\int o_M(d)\,\xi(dd).
\]
The ordinary error under $\mu$ is at most $\eta_\mu(M)$.
Stop as soon as either component stops and return that output.
On completed counterfactual streams, the physical work is at most
\[
 2\min\{T_{\rm bank},T_{\rm ordinary}\}\le2M.
\]
A wrong output implies either a wrong bank stop or an ordinary error.
When $\sigma=0$, omit the bank and use the ordinary classifier alone.

On a $\sigma$-almost-every gap $d$ in band $j$, the matching negative
reference has drift at least $4d^2/3$ and log overshoot below one.
Its expected crossing time is at most
\begin{equation}\label{eq:split-negative-work}
 \frac3{4d^2}[1+\log(1/\epsilon)+\log(1/\pi_j)].
\end{equation}
On a positive gap, the event $\widehat p_n-1/2\ge d/2$ gives
\begin{equation}\label{eq:split-null-laplace}
 \frac{P_{-,\mu,n}}{P_{{\rm f},n}}
 \le\int e^{-2nx^2}\,\mu(dx)
 \le z+\frac A{2en}+\frac H{2\sqrt{en}}.
\end{equation}
The last two bounds use, respectively,
$\sup_xx^2e^{-2nx^2}=1/(2en)$ and
$\sup_xxe^{-2nx^2}=1/(2\sqrt{en})$.
For $n\ge d^{-2}$ the denominator is at most $z+Ad^2+Hd$,
while the matching positive reference contributes at least
$\pi_je^{nd^2/3}$ relative to fair. Equation~\eqref{eq:uf-bank-positive} therefore bounds positive expected bank work by
\begin{equation}\label{eq:split-positive-work}
 \frac6{d^2}\left[
 1+\log(1/\epsilon)+\log(1/\pi_j)
 +\log_+\frac{z+Ad^2+Hd}{\rho}\right].
\end{equation}
These arguments prove finite conditional bank work on $\sigma$.
On $\tau+\zeta$, the ordinary cap guarantees termination of the race.

Since $z/\rho\le1$,
\[
 \log_+\frac{z+Ad^2+Hd}{\rho}
 \le\log3+\log_+(Ad^2/\rho)+\log_+(Hd/\rho).
\]
Using $\sup_{x\ge1}\log x/x=1/e$ and
$\sup_{x\ge1}\log x/x^2=1/(2e)$ gives
\begin{equation}\label{eq:split-surcharge}
 \int q(d)d^{-2}\log_+\frac{z+Ad^2+Hd}{\rho}\,\sigma(dd)
 \le(\log3+1/e)A+\frac{H^2}{2e\rho}.
\end{equation}
Here $\int q\,d\sigma\le\rho$ and $q\le1$ follow from the original
global rarity class; no separate conditional rarity assumption is used.
Integrating \eqref{eq:split-negative-work} and
\eqref{eq:split-positive-work} therefore bounds the bank work on $\sigma$ by
\begin{equation}\label{eq:split-bank-work}
 6\Gamma_\sigma+6A\log(1/\epsilon)
       +6(\log3+1/e)A+\frac{3H^2}{e\rho}.
\end{equation}

\subsection{A resolved-law cap schedule}

Let $\chi=\sigma+\tau$ and set
\[
 \delta_\ell=\rho16^{-(\ell+1)},\qquad
 u_\ell=\delta_\ell/4,\qquad \epsilon_\ell=u_\ell/\rho.
\]
Choose $M_\ell$ as the least power of four, at least four, such that
$\eta_\chi(M_\ell)\le\delta_\ell/8$.
Every chosen cap has a predecessor because $\eta_\chi(4)=1-z$.
The inequality $\eta_\chi(4M)\ge\eta_\chi(M)/16$ implies
\begin{equation}\label{eq:split-cap-selection}
 \delta_\ell/128<\eta_\chi(M_\ell)\le\delta_\ell/8.
\end{equation}
Thus the selected powers of four are distinct and increasing.

If $z>0$, let $J$ be the largest integer with $\delta_J\ge8z$.
Then
\begin{equation}\label{eq:split-terminal-level}
 8z\le\delta_J<128z.
\end{equation}
For $\ell\le J$,
$\eta_\mu(M_\ell)\le\eta_\chi(M_\ell)+z\le\delta_\ell/4=u_\ell$.
The raced classifier has total error at most $3u_\ell\le\delta_\ell$
and false-positive probability at most $2u_\ell\le\delta_\ell$.
When $z=0$, these conclusions hold at every level.

Since $\log(1/\epsilon_\ell)=\log64+\ell\log16$ and
$A\le\Gamma_\sigma$, \eqref{eq:split-bank-work} is at most
$40(\ell+1)\Gamma_\sigma+2H^2/\rho$.
A uniform raced-work certificate is consequently
\begin{equation}\label{eq:split-level-work}
 C_\ell=80(\ell+1)\Gamma_\sigma+4H^2/\rho
                  +2O_\tau(M_\ell)+2zM_\ell.
\end{equation}
It remains a valid, possibly loose, bound when $\sigma=0$ and the bank is
omitted.

The binary level-zero output has second moment at most
$\rho+\delta_0=17\delta_0$. An adjacent difference, run on successive
unused blocks of the same coin, has second moment at most
$\delta_\ell+\delta_{\ell-1}=17\delta_\ell$.
Set $v_\ell=17\delta_\ell$, $k_0=C_0$,
$k_\ell=C_\ell+C_{\ell-1}$ for $\ell\ge1$, and
$S_J=\sum_{\ell\le J}\sqrt{v_\ell k_\ell}$.
Independent-Poisson allocation across these corrections
(Lemma~\ref{lem:uf-poisson})
gives RMS at most
\begin{equation}\label{eq:split-poisson}
 \sqrt{\delta_J^2+S_J^2/B}.
\end{equation}
Every such finite-level procedure has finite deterministic local caps.
In the finite-law theorem the reference bank is finite; its stopping
events are measurable and its positive reference is a finite probability
mixture. No countable-support extension is needed here.

\subsection{Controlling mixed work without an extra logarithm}

For every finite $\tau$, the following bounds hold:
\begin{equation}\label{eq:split-elementary-kernels}
 O_\tau(M)\le21H\sqrt M,\qquad
 \eta_\sigma(M)\le A/M,\qquad
 \eta_\tau(M)\le H/\sqrt M.
\end{equation}
The latter two are pointwise inequalities before integration.
For the first, put $x=\sqrt M\,d$. The ratio of $o_M(d)$ to
$\sqrt M/d$ is at most one for $x\le1$, and at most
$17(1+2\log x)/x\le34/\sqrt e<21$ for $x\ge1$.
Lemma~\ref{lem:uf-kernel} supplies the homogeneous ordinary kernel bound:
\begin{equation}\label{eq:split-ordinary-kernel}
 \sum_{k\ge1}\sqrt{\eta_\tau(4^k)O_\tau(4^k)}\le 208H.
\end{equation}
Its proof applies without normalization of $\tau$: dyadic binning yields
the summable kernels $2^{k-j}$ for $j\ge k$, and
$2^{3(j-k)}$, $2^{j-k}(1+k-j)$ for $j<k$.
Apply $\sqrt{xy}\le(x+y)/2$ and sum these kernels.

Split levels according to which term of
$\eta_\chi=\eta_\sigma+\eta_\tau$ is at least half the total.
On levels where $\eta_\tau$ dominates,
$\delta_\ell\le256\eta_\tau(M_\ell)$, so distinctness of the caps and
\eqref{eq:split-ordinary-kernel} give a total
$\sum\sqrt{\delta_\ell O_\tau(M_\ell)}\le3328H$.
On the other levels,
\[
 \delta_\ell\le256\eta_\sigma(M_\ell)\le256A/M_\ell,
 \qquad M_\ell\le256A/\delta_\ell.
\]
Therefore
\[
 \sqrt{\delta_\ell O_\tau(M_\ell)}
 \le4\sqrt{21}\,H^{1/2}A^{1/4}\delta_\ell^{1/4}.
\]
The sum of $\delta_\ell^{1/4}$ over all levels is $\rho^{1/4}$.
The arithmetic--geometric mean inequality proves
\begin{equation}\label{eq:split-mixed-sum}
 \sum_{\ell\le J}\sqrt{\delta_\ell O_\tau(M_\ell)}
 \le C_{\rm mix}[H+\sqrt{\rho A}].
\end{equation}

The work on the unresolved mass is controlled separately.
On levels where $\eta_\sigma$ dominates,
\begin{equation}\label{eq:split-unresolved-sigma}
 \sum\sqrt{z\delta_\ell M_\ell}
 \le16(J+1)\sqrt{zA}\le\sqrt{2\rho A}.
\end{equation}
The last inequality uses $z\le\delta_J/8$ and
$(J+1)4^{-(J+1)}\le1/4$.
On levels where $\eta_\tau$ dominates,
$\sqrt{M_\ell}\le256H/\delta_\ell$, so
\begin{equation}\label{eq:split-unresolved-tau}
 \sum\sqrt{z\delta_\ell M_\ell}
 \le256H\sqrt z\sum_{\ell=0}^J\delta_\ell^{-1/2}
 \le\frac{1024}{3}H\sqrt{z/\delta_J}
 \le\frac{1024}{3\sqrt8}H.
\end{equation}
When $z=0$ these contributions vanish.

\begin{proof}[Completion of Theorem~\ref{thm:direct-upper}]
The shifted adjacent-level sum has a factor $1/4$, so
\[
 S_J\le\frac{5\sqrt{17}}4
                \sum_{\ell=0}^J\sqrt{\delta_\ell C_\ell}.
\]
Use \eqref{eq:split-level-work}, \eqref{eq:split-mixed-sum},
\eqref{eq:split-unresolved-sigma}, \eqref{eq:split-unresolved-tau},
and $A\le\Gamma_\sigma$. The elementary sums
$\sum_{\ell\ge0}4^{-\ell}\sqrt{\ell+1}\le16/9$ and
$\sum_{\ell\ge0}4^{-\ell}=4/3$ then give
\[
 S_J\le C_{\rm sum}[\sqrt{\rho\Gamma_\sigma}+H].
\]
One valid explicit choice is $C_{\rm sum}=30000$. Indeed, writing
$R=\sqrt{\rho\Gamma_\sigma}$, the component bounds above give
\[
 \sum_{\ell\le J}\sqrt{\delta_\ell C_\ell}
 \le21R+\frac{15535}{3}H.
\]
Here the core term is at most $4R$, the fixed tail surcharge is $2H/3$,
the ordinary mixed term is at most $15R+5007H$, and the unresolved
term is at most $2R+512H/3$. Use $\sqrt{21}<5$, $\sqrt2<3/2$ and
$5\sqrt{17}/4<165/32$ to obtain
\[
 S_J\le\frac{3465}{32}R+\frac{854425}{32}H
       \le30000(R+H).
\]
For $z>0$, \eqref{eq:split-terminal-level} and
\eqref{eq:split-poisson} imply
$R(B)\le128z+C_{\rm sum}[\sqrt{\rho\Gamma_\sigma}+H]/\sqrt B$.
For $z=0$, put $D=\sqrt{\rho\Gamma_\sigma}+H>0$ and choose the least
finite $J\ge0$ such that $\delta_J\le D/\sqrt B$. The bound on $S_J$
holds at this finite terminal level. Equation~\eqref{eq:split-poisson}
therefore gives RMS at most
\[
 \sqrt{1+30000^2}\,D/\sqrt B<30001D/\sqrt B.
\]
Thus this branch has an actual finite-level estimator; no limiting
or infinite-level construction is needed.
The constant estimator covers $z>\rho/128$.
For the fixed split, $z+hH=T_h$ and $\sigma=\kappa$.
The preceding bound and the constant estimator therefore give
$R(B)\le C\min\{\rho,T_h+h\sqrt{\rho\Gamma_\kappa}\}=C\Q(B)$.
The preceding constants permit $C=30001$ for this upper bound; no
optimality of that constant is claimed. The number of levels and every
local cap are finite, while the total Poisson replication count is
unbounded but almost surely finite and has the stated expected budget.
Since the law is finite, the reference bank is finite throughout.
\end{proof}

\section{Completion of the main theorem}\label{sec:assembly-contract}
The same finite known law $\mu$, rarity bound $\rho$, parameter class
$\mathcal Q_\rho$, persistent latent-pair experiment, target $\theta(q)$,
RMS loss, and uniform expected observation budget are used on both sides.
For every fixed $(\mu,\rho,B)$, the upper proof gives one observable
algorithm whose work and risk bounds hold for every $q\in\mathcal Q_\rho$.
Its design may depend on $(\mu,\rho,B)$ but not on $q$ or a genuine
coin's latent state. The lower proof applies to every such algorithm,
including nonlinear outputs and arbitrary interleaving.
There is no claim of a single procedure simultaneously attaining every
budget, or of a deterministic total observation cap.

\subsection{Upper-bound branches and actual attainment}
Put $D=\sqrt{\rho\Gamma_\kappa}+H$ for the split in
Section~\ref{sec:direct-upper}. The following exhaustive branches are
selected using only known design parameters:
\begin{enumerate}
\item If $\Q=\rho$, output $\rho/2$ at zero cost.
\item If $\Q<\rho$ and $z>\rho/128$, the same constant has risk
$\rho/2<64z\le64\Q$.
\item If $\Q<\rho$ and $0<z\le\rho/128$, then
$\delta_0=\rho/16\ge8z$, so the largest $J$ with $\delta_J\ge8z$
exists and is finite. The error is at most
$128z+30000hD\le30000\Q$.
\item If $\Q<\rho$ and $z=0$, then $D>0$ since the core has positive
mass. The least $J$ with $\delta_J\le hD$ exists and is finite;
the error is less than $30001hD=30001\Q$.
\end{enumerate}
In the last two branches every local cap exists: the finite resolved law
has positive gaps and $\eta_\chi(M)\to0$. Moreover
$\eta_\chi(4)=1-z>\delta_\ell/8$, including at the boundary $z=\rho/128$.
Poisson means are positive and finite. Independently drawn counts and
the uniform per-replication certificates give, for every $q$,
\[
 \E_qN\le\sum_{\ell=0}^J n_\ell k_\ell=B.
\]
The counts are almost surely finite and all local caps are finite.
Thus an actual admissible finite-level procedure exists at every $B>0$;
no attainment of a minimax infimum or infinite-level limit is assumed.
It follows that $R(B)\le30001\Q(B)$.

\subsection{Lower-bound constants and separate hard parameters}
Put $M_B=\min\{\rho,h\sqrt{\rho F}\}$ and
\[
 c_{\rm ann}=\frac{e^{-28}}{2000000},\qquad
 k=\frac{\sqrt{c_{\rm ann}}}{128\cdot1024\cdot3584}.
\]
In the packing branch, additive localization gives
\[
 R(B)\ge\frac1{128}\sum_i
      \min\{\theta_i/64,\sqrt{c_{\rm ann}}\,h\sqrt{\theta_iG_i}\}
 \ge k M_B.
\]
The last step uses $\sqrt{c_{\rm ann}}<1/64$, the proportional rarity
allocation, and $\sum_iG_i\ge F/3584$.
In the high-band branch $F\le16C_H$ and
$C_H=4096(1+\log6)<16384$. Hence the Bernoulli-label floor yields
$R(B)\ge M_B/16384\ge kM_B$ there as well.

If $h\ge t$, the tail bound gives $R\ge\Q/32$.
If $h<t$ and $\lambda\ge t/4$, it gives $R\ge\Q/288$.
Otherwise both $R\ge T_h/32$ and $R\ge kM_B$ hold. Since
$F\le\Gamma_\kappa\le4F$,
\[
 \Q\le T_h+2M_B\le(32+2/k)R,
 \qquad R\ge\frac{k}{32k+2}\Q\ge\frac{k}{3}\Q.
\]
Here $32k\le1$. The other branches have larger constants, so this proves
one lower constant uniformly over all laws, rarities and budgets.
The tail and core obstructions need not be witnessed by the same
parameter: a worst-case risk exceeds each obstruction, hence their
maximum. No exchange of a supremum and a sum is being used.

The explicit bound in Theorem~\ref{thm:direct-main} is therefore
\[
             10^{-20}\Q(B)\le R(B)\le30001\Q(B).
\]
Indeed $k/3=e^{-14}/[\sqrt{2000000}\,(3\cdot128\cdot1024\cdot3584)]$,
and $e<3$, $\sqrt{2000000}<1500$ and
\[
 3^{14}\,1500\,(3\cdot128\cdot1024\cdot3584)<10^{20}
\]
certify this deliberately loose lower constant.

\subsection{Lower-bound quantifiers and budget transport}
The half-prior components are subsets of the full rarity class after
the fractional mixture is integrated. Each induced endpoint and its
population midpoint is admissible, so the original uniform budget
applies to all three. Auxiliary oracle types strengthen only the lower
experiment; they are not observations available to the upper estimator.
The original algorithm is simulated using persistent private labels,
and genuine oracle observations never exceed its original work.

At each fixed budget, the common-pair result selects a pair simultaneously
for all local policies and deterministic caps. That pair is held fixed
while the cap tends to infinity. Population midpoint work is bounded by
$B$, and the overflow probabilities at the two endpoints sum to at most
$2B/L$ at cap $L$. Neither a null-only nor an endpoint-only budget is
substituted for the uniform assumption. The population midpoint is the
experiment with an independent midpoint latent pair for every fresh
coin, not the mixture of the two full population endpoint transcripts.

This completes the proof of Theorem~\ref{thm:direct-main}.

\section{Consequences}\label{sec:consequences}
Write $c_0=10^{-20}$ and $C_0=30001$ for the constants in
Theorem~\ref{thm:direct-main}.

\begin{corollary}[Mean-square form]\label{cor:mse}
For every law, rarity and budget in Theorem~\ref{thm:direct-main},
\[
 c_0^2\min\{\rho^2,T_h^2+\rho\Gamma_\kappa/B\}
 \le R(B)^2
 \le2C_0^2\min\{\rho^2,T_h^2+\rho\Gamma_\kappa/B\}.
\]
If $B\ge d_1^{-2}$ and $J_\nu=\int d^{-1}\nu(dd)$, then
\[
 c_0^2\min\{\rho^2,(J_\nu^2+\rho\Gamma_\kappa)/B\}
 \le R(B)^2
 \le2C_0^2\min\{\rho^2,(J_\nu^2+\rho\Gamma_\kappa)/B\}.
\]
\end{corollary}
\begin{proof}
For $x,y\ge0$, $x^2+y^2\le(x+y)^2\le2(x^2+y^2)$.
Taking the minimum with $\rho^2$ preserves these comparisons,
with the same factor two. Apply this to $x=T_h$ and
$y=h\sqrt{\rho\Gamma_\kappa}$ and square
\eqref{eq:main-bounds}. If $B\ge d_1^{-2}$, then $h\le d_i$
at every atom, so $T_h=hJ_\nu$ exactly, including at equality.
\end{proof}

\begin{corollary}[One-gap specialization]\label{cor:one-gap}
For $\mu=\delta_d$, $0<d\le1/2$,
\[
 R_{\delta_d,\rho}(B)\asymp
       \min\{\rho,\sqrt{\rho}/(d\sqrt B)\}.
\]
\end{corollary}
\begin{proof}
Here $\nu=\rho\delta_d$, $\kappa=(1-\rho)\delta_d$, and
$\Gamma_\kappa=(1-\rho)/d^2$. Put $x=h/d$. If $x\ge1$,
then $\Q(B)=\rho$. If $x<1$, the sum inside the minimum
defining $\Q$ is $x[\rho+\sqrt{\rho(1-\rho)}]$.
Since $0<\rho\le1/4$, the bracket lies between
$(\sqrt3/2)\sqrt\rho$ and $(3/2)\sqrt\rho$.
Taking the minimum with $\rho$ and applying the theorem proves
the comparison with universal constants.
\end{proof}

For a general finite law the core entropy vanishes only in the
normalized entropy term when a single band is occupied; its base
cost $A^*$ is always positive. The quantile split can place part
of one atom in both $\nu$ and $\kappa$. These cases are included
in the theorem without any generic-position assumption.

The formulas concern a uniform expected observation budget.
They do not supply a deterministic total cap, a high-confidence
guarantee, or an estimator that learns an unknown magnitude law.
The finite support and strictly positive gaps ensure that all
caps used in the upper construction exist. Extending these
assumptions requires additional arguments.

\clearpage
\appendix
\section{Observable classification and finite Poisson allocation}
\label{sec:uf-foundations}
This appendix proves the single-coin classification bounds and the
finite multilevel estimator used in Section~\ref{sec:direct-upper}.
Throughout this appendix, a false-positive probability means the joint
event $\Pp_q(Z=1,s=0)$, with the binary sign $s=\mathbf1_{\{S=+1\}}$, not the conditional
probability $\Pp_q(Z=1\mid s=0)$. All likelihoods
refer to the actual binary observations. A reference distribution evaluates
that prefix and does not reveal or redraw the latent state of a coin.

\subsection{A capped classifier that does not know the gap}

For a dyadic integer $M\ge1$, inspect prefix lengths
$1,2,4,\ldots,M$. At a prefix $n<M$, stop if
\[
 \left|\widehat p_n-\frac12\right|>
 t_{n,M}:=\sqrt{\frac{1+\log(M/n)}{n}},
\]
and output the side of $1/2$ containing $\widehat p_n$. At $M$, output
$\1_{\{\widehat p_M>1/2\}}$, with a negative output on ties.
Denote the output and work by $(Z_M^{\rm ord},T_M^{\rm ord})$.

\begin{lemma}[Ordinary conditional error and work]\label{lem:uf-ordinary}
For every $0<d\le1/2$ and either latent sign,
\begin{align}
 \Pp_{d,s}(Z_M^{\rm ord}\ne s)
 &\le \min\{1,(Md^2)^{-2}\},\label{eq:uf-ordinary-error}\\
 \E_{d,s}T_M^{\rm ord}
 &\le o_M(d):=\min\{M,17d^{-2}[1+\log_+(Md^2)]\}.
 \label{eq:uf-ordinary-work}
\end{align}
\end{lemma}

\begin{proof}
For independent Bernoulli observations of mean $p$, Hoeffding's
inequality gives
$\Pp(\widehat p_n-p\ge t)\le e^{-2nt^2}$ and the analogous lower-tail
bound \cite{hoeffding}. A wrong early output at a fixed inspected prefix
requires a deviation of at least $d+t_{n,M}$. Applying a union bound on
the underlying full tape, without conditioning on the stopping event,
therefore gives
\[
 \Pp_{d,s}(Z_M^{\rm ord}\ne s)
 \le e^{-2Md^2}
       +e^{-2}\sum_{\substack{n<M\\n\ {\rm dyadic}}}
                    (n/M)^2e^{-2nd^2}.
\]
For any dyadic sequence of positive $u$ values, the intervals
$[u/2,u]$ have disjoint interiors and
\[
 \int_{u/2}^{u}y e^{-2y}\,dy
 \ge\frac38u^2e^{-2u}.
\]
Consequently $\sum_u u^2e^{-2u}\le2/3$.
Writing $x=Md^2\ge1$ and using $x^2e^{-2x}\le e^{-2}$, the preceding
error bound is at most $(5/3)e^{-2}x^{-2}<x^{-2}$. For $x<1$, the
probability bound by one suffices.

Since the stopping time takes dyadic values,
\[
 \E_{d,s}T_M^{\rm ord}
 =1+\sum_{\substack{n<M\\n\ {\rm dyadic}}}
                      n\,\Pp_{d,s}(T_M^{\rm ord}>n).
\]
If $d>2t_{n,M}$, survival implies a deviation from the true mean of
at least $d/2$. Thus
\[
 \E_{d,s}T_M^{\rm ord}
 \le1+\sum_{\substack{n<M\\n\ {\rm dyadic}}}
       n\left[\1_{\{d\le2t_{n,M}\}}+2e^{-nd^2/2}\right].
\]
For $x=Md^2\ge1$, the first indicator implies
$nd^2\le4(1+\log x)$. Indeed, this is immediate if $nd^2<1$; otherwise
$\log(M/n)=\log(x/(nd^2))\le\log x$.
The corresponding dyadic sum is at most
$8d^{-2}(1+\log x)$.
Also
\[
 \int_{u/2}^{u}e^{-y/2}\,dy\ge\frac12u e^{-u/2}
 \quad\Longrightarrow\quad
 \sum_u u e^{-u/2}\le4.
\]
The exponential part of the work bound is therefore at most $8d^{-2}$.
Since $1\le d^{-2}$, total expected work is at most
$17d^{-2}(1+\log x)$. Finally $T_M^{\rm ord}\le M$ deterministically;
when $x<1$, the displayed minimum equals $M$.
\end{proof}

For any finite measure $\xi$ on $(0,1/2]$, write
\[
 \eta_\xi(M)=\int\min\{1,(Md^2)^{-2}\}\,\xi(dd),
 \qquad
 O_\xi(M)=\int o_M(d)\,\xi(dd).
\]
Both are homogeneous in $\xi$. The first is nonincreasing in $M$ and
the second is nondecreasing. Pointwise comparison also gives
$\eta_\xi(4M)\ge\eta_\xi(M)/16$.

\begin{lemma}[A homogeneous dyadic kernel estimate]\label{lem:uf-kernel}
For every finite measure $\xi$ and integer $K\ge1$,
\[
 \sum_{k=1}^{K}\sqrt{\eta_\xi(4^k)O_\xi(4^k)}
 \le208\int\min\{2^K,d^{-1}\}\,\xi(dd).
\]
In particular,
\begin{equation}\label{eq:uf-kernel-infinite}
 \sum_{k\ge1}\sqrt{\eta_\xi(4^k)O_\xi(4^k)}
 \le208\int d^{-1}\,\xi(dd),
\end{equation}
with extended-value interpretation.
\end{lemma}

\begin{proof}
Use bands $(2^{-j-1},2^{-j}]$, $j\ge1$, with masses $w_j$, and put
$c_j=2^jw_j$. Direct substitution into the two kernels yields
\[
 \eta_\xi(4^k)\le2^{-k}
 \left[\sum_{j\ge k}c_j2^{k-j}
              +16\sum_{j<k}c_j2^{3(j-k)}\right]=2^{-k}A_k,
\]
\[
 O_\xi(4^k)\le2^k
 \left[\sum_{j\ge k}c_j2^{k-j}
        +136\sum_{j<k}c_j2^{j-k}(1+k-j)\right]=2^kD_k.
\]
For the second display on $j<k$, use
$d^{-2}\le4\,2^{2j}$ and
$1+\log(4^kd^2)\le2(1+k-j)$.
On $j\ge k$, the deterministic cap gives the first sum exactly as
an upper bound.

Apply $\sqrt{A_kD_k}\le(A_k+D_k)/2$ and sum over $k\le K$.
For a fixed $j\le K$, the three relevant geometric sums obey
\[
 \sum_{k\le j}2^{k-j}\le2,\qquad
 \sum_{r\ge1}2^{-3r}=\frac17,\qquad
 \sum_{r\ge1}2^{-r}(1+r)=3.
\]
The coefficient of $c_j$ is thus at most
$2+8/7+204<208$. For $j>K$, only the first kernels occur, with
coefficient at most $2\,2^{K-j}$. Hence the result is at most
\[
 208\sum_{j\le K}2^jw_j+2\,2^K\sum_{j>K}w_j.
\]
For $j\le K$, $\min\{2^K,d^{-1}\}\ge2^j$ on the band; for $j>K$
it equals $2^K$. This proves the finite bound. Increasing $K$ and
using monotone convergence gives \eqref{eq:uf-kernel-infinite}.
All rearrangements involve nonnegative terms, so normalization and
moment assumptions are unnecessary.
\end{proof}

\subsection{A weighted bank compared with the full null mixture}

For a binary prefix $x_{1:n}$, let $P_{p,n}(x_{1:n})$ be its probability
under iid $\operatorname{Ber}(p)$ observations, and let
$P_{{\rm f},n}=P_{1/2,n}$. Choose any nonempty finite or countable set
of canonical bands
\[
 I_j=(b_j/2,b_j],\qquad b_j=2^{-j-1},
\]
and positive weights $\pi_j$ summing to one on these bands.
Set $r_j^-=b_j/2$ and $r_j^+=b_j/4$, and define
\[
 L_{j,n}^-=\frac{P_{1/2-r_j^-,n}}{P_{{\rm f},n}},
 \qquad
 M_n=\frac{\sum_j\pi_jP_{1/2+r_j^+,n}}
           {P_{-,\mu,n}},
 \qquad
 P_{-,\mu,n}=\int P_{1/2-d,n}\,\mu(dd).
\]
For $0<\epsilon<1$, stop negative when any
$L_{j,n}^-\ge1/(\epsilon\pi_j)$, and stop positive when
$M_n\ge1/(\rho\epsilon)$. Give negative crossings priority.
A zero denominator triggers a positive stop. Denote the bank stopping
time by $T_{\rm bank}$; it may be infinite outside its selected bands.

\begin{lemma}[Full-mixture likelihood bank]\label{lem:uf-mixture-bank}
The probability of a wrong bank stop is at most $2\rho\epsilon$
uniformly over $q\in\mathcal Q_\rho$, and its false-positive probability
is at most $\rho\epsilon$. These statements concern the event that the
bank stops and do not assume termination at every magnitude.

For $d\in I_j$ in a selected band,
\begin{equation}\label{eq:uf-bank-negative}
 \E_{d,-}T_{\rm bank}
 \le\frac3{4d^2}
       [1+\log(1/\epsilon)+\log(1/\pi_j)].
\end{equation}
Suppose a finite number $V(d)>0$ satisfies
\[
 \beta_\mu(n):=\int e^{-2nx^2}\,\mu(dx)\le V(d)
 \qquad\text{for every integer }n\ge d^{-2}.
\]
Then
\begin{equation}\label{eq:uf-bank-positive}
 \E_{d,+}T_{\rm bank}
 \le\frac6{d^2}
 \left[1+\log(1/\epsilon)+\log(1/\pi_j)
                   +\log_+\frac{V(d)}{\rho}\right].
\end{equation}
In particular the bank has finite conditional work on every selected
band whenever such an envelope is supplied.
\end{lemma}

\begin{proof}
A nonnegative supermartingale $Y_n$ with $Y_0=1$ has
$\Pp(\sup_nY_n\ge a)\le1/a$: stopping at the first crossing or a fixed
horizon gives $1\ge\E Y_{T\wedge m}\ge a\,\Pp(T\le m)$, and increasing
the horizon proves the claim. This elementary maximal inequality also
works when the process is a supermartingale only on supported prefixes.

Under a positive coin of gap $d$, the next multiplier of $L_{j,n}^-$
has mean $1-4r_j^-d\le1$. Its negative-crossing probability is at most
$\epsilon\pi_j$. A countable union bound gives conditional
negative-crossing probability at most $\epsilon$, so its positive-sign
contribution under any admissible prior is at most $\rho\epsilon$.

Under the full all-negative mixture, $M_n$ is a nonnegative
supermartingale on supported prefixes. If a prefix $x$ has positive
null probability, conditional expectation sums only children $xu$
with positive null probability; their numerator probabilities sum to
at most the numerator at $x$. This proves the supermartingale
inequality even when $\mu$ has mass at $d=1/2$. A zero-denominator prefix
has null probability zero. Thus a positive crossing has null-mixture
probability at most $\rho\epsilon$. Its actual negative-sign
contribution is no larger, since it multiplies the conditional
crossing probabilities by $1-q(d)\le1$. The two error claims follow.

For negative work, fix a selected-band gap $d\in I_j$ and
$r=r_j^-\in[d/2,d]$, with $r\le1/4$. The matching log likelihood has
mean increment
\[
 g(d,r)=\frac12\log(1-4r^2)
          +d\log\frac{1+2r}{1-2r}
 \ge-\frac83r^2+4dr\ge\frac43d^2.
\]
Here $\log(1-x)\ge-x/(1-x)$ and
$\log((1+x)/(1-x))\ge2x$ establish the first lower bound; the
concave quadratic takes the value $4d^2/3$ at both endpoints
$r=d/2,d$. Every upward log increment is at most $\log(3/2)<1$.
For the first correct-reference crossing $\tau$, the stopped-sum
identity at $\tau\wedge m$ gives
\[
 g(d,r)\E(\tau\wedge m)
 =\E\log L^-_{j,\tau\wedge m}
 \le1+\log(1/(\epsilon\pi_j)).
\]
Monotone convergence proves integrability of $\tau$ and the stated
work bound, since $T_{\rm bank}\le\tau$. No unbounded stopping identity
is invoked before integrability is established.

For positive work, let
$E_n=\{\widehat p_n-1/2\ge d/2\}$. On this event, for $x<1/2$,
\[
 \log\frac{P_{1/2-x,n}}{P_{{\rm f},n}}
 =\frac n2\log(1-4x^2)
   -n(\widehat p_n-\tfrac12)
                 \log\frac{1+2x}{1-2x}
 \le-2nx^2.
\]
For $x=1/2$ the likelihood is zero, since $E_n$ entails at least one
head. Therefore $P_{-,\mu,n}/P_{{\rm f},n}\le\beta_\mu(n)$.
The matching positive reference $r=r_j^+\in[d/4,d/2]$ satisfies,
on the same event,
\[
 \log\frac{P_{1/2+r,n}}{P_{{\rm f},n}}
 \ge n\left[-\frac83r^2+2dr\right]\ge nd^2/3.
\]
Set
\[
 U=1+\log(1/\epsilon)+\log(1/\pi_j)
                  +\log_+(V(d)/\rho),\qquad
 n_0=\left\lceil3d^{-2}U\right\rceil.
\]
For every $n\ge n_0$, $E_n$ implies
$M_n\ge\pi_je^{nd^2/3}/V(d)\ge1/(\rho\epsilon)$.
Thus $T_{\rm bank}>n$ implies $E_n^c$, whose positive-coin probability
is at most $e^{-nd^2/2}$. Summing the survival probabilities yields
\[
 \E_{d,+}T_{\rm bank}
 \le3d^{-2}U+2+2d^{-2}\le6d^{-2}U.
\]
The last inequality uses $U\ge1$ and $d^2\le1/4$.

For a countable bank, the positive mixture is a consistent probability
law and the negative stopping event is a measurable countable union.
At any fixed time only finitely many negative references can cross:
$L^-_{j,n}\le(3/2)^n$ requires
$\pi_j\ge1/[\epsilon(3/2)^n]$ for a crossing. All preceding arguments
therefore apply to finite and countable banks alike.
\end{proof}

\subsection{Racing procedures and retaining a genuine unit}

\begin{lemma}[Observable race]\label{lem:uf-race}
Suppose two classifiers run on the odd and even subsequences of one
unit's observation tape, with independent auxiliary seeds. Stop the
physical procedure when either classifier stops and use its output.
On completed counterfactual subsequences its work satisfies
\[
 T_{\rm race}\le2\min\{T_1,T_2\}.
\]
Its error event is contained in the union of the component wrong-stop
events. If either component has a finite deterministic cap, the race
terminates with a finite deterministic cap.
\end{lemma}

\begin{proof}
By physical time $2n$, each component has received $n$ observations.
If either component would stop after $n$ of its own observations, the
race stops no later than $2n$. This proves the pathwise work bound,
including an infinite stopping time for the other component.
The selected output is wrong only if the component that produces it
makes a wrong stop. Conditional on the latent $(D,S)$, both
subsequences are iid Bernoulli tapes of the correct mean, so all
component error and work bounds apply. The policy uses only the
physical observed prefix and its seeds.
\end{proof}

\subsection{Finite multilevel allocation at every positive budget}

\begin{lemma}[Independent Poisson allocation]\label{lem:uf-poisson}
Let $Z_0,\ldots,Z_J\in\{0,1\}$ be observable single-unit classifiers,
each with finite deterministic cap. Suppose uniformly over
$q\in\mathcal Q_\rho$ that their error probabilities are at most
$e_\ell$, their false-positive probabilities are at most $f_\ell$,
and their expected work is at most $C_\ell>0$.
Put
\[
 v_0=\rho+f_0,\qquad
 v_\ell=e_\ell+e_{\ell-1}\quad(\ell\ge1),\qquad
 k_0=C_0,\qquad k_\ell=C_\ell+C_{\ell-1}\quad(\ell\ge1).
\]
Assume $v_\ell>0$ and set $A_J=\sum_{\ell=0}^J\sqrt{v_\ell k_\ell}$.
For every $B>0$ there is an observable population estimator of
uniform expected work at most $B$ and worst RMS error at most
\begin{equation}\label{eq:uf-poisson-risk}
 \sqrt{e_J^2+A_J^2/B}.
\end{equation}
\end{lemma}

\begin{proof}
A level-zero correction uses a fresh genuine unit and returns
$\Delta_0=Z_0$. For $\ell\ge1$, a correction draws one fresh latent unit,
runs $Z_\ell$ on its observations, then runs $Z_{\ell-1}$ with a fresh
seed on the next unused observations of the same unit, and returns
$\Delta_\ell=Z_\ell-Z_{\ell-1}$.

Conditional on the latent state, the unused suffix after a finite
stopping time is an iid tape with the same Bernoulli mean: for each
possible stopped prefix, independence of its subsequent coordinates
gives the claim, and summing over the stopped prefixes preserves it.
The second classifier therefore has its original marginal distribution
under the same latent state. No new magnitude or sign is drawn between
the two classifiers. Their disagreement implies that at least one
misclassifies that common sign, and hence
\[
 \E_q\Delta_\ell^2\le v_\ell,\qquad
 \E_q W_\ell\le k_\ell,\qquad
 \sum_{\ell=0}^J\E_q\Delta_\ell=\E_qZ_J.
\]
At level zero use
$\Pp_q(Z_0=1)\le\Pp_q(s=1)+f_0\le\rho+f_0$.

Before drawing the genuine units and their observation tapes, draw
the Poisson replication counts independently across levels, with means
\[
 n_\ell=\frac B{A_J}\sqrt{\frac{v_\ell}{k_\ell}}.
\]
For each replication use an independent genuine unit, and return
\[
 \widehat\theta
 =\sum_{\ell=0}^J\frac1{n_\ell}
       \sum_{i=1}^{N_\ell}\Delta_{\ell,i},
 \qquad N_\ell\sim\operatorname{Poisson}(n_\ell).
\]
Its mean is $\E_qZ_J$, so its absolute bias is at most $e_J$.
The expected work is at most
$\sum_\ell n_\ell k_\ell=B$.

For independent marks $\Delta$ and a Poisson count of mean $n$,
conditioning on the count gives
$\Var(\sum_{i=1}^{N}\Delta_i)=n\E\Delta^2$.
Consequently the independent level variances sum to at most
\[
 \sum_{\ell=0}^J\frac{v_\ell}{n_\ell}=\frac{A_J^2}{B}.
\]
The divisor is the deterministic mean $n_\ell$, including when the
realized count is zero; it is not the random count $N_\ell$. Count
fluctuations contribute the squared mark mean, which is why the second
moment rather than just the mark variance appears. The mark and its
observation cost may be dependent within a replication; the count is
independent of the entire sequence of these paired draws.

For any positive allocation $x_\ell$ with $\sum_\ell k_\ell x_\ell\le B$,
Cauchy--Schwarz gives
\[
 \left(\sum_\ell\sqrt{v_\ell k_\ell}\right)^2
 \le\left(\sum_\ell v_\ell/x_\ell\right)
       \left(\sum_\ell k_\ell x_\ell\right).
\]
Thus the displayed allocation minimizes this second-moment certificate
under its expected-work constraint. This is not a deterministic work cap.
Bias squared plus variance proves \eqref{eq:uf-poisson-risk}.
All counts are almost surely finite, all component caps are finite,
and all first and subsequent observations are included in $W_\ell$.
\end{proof}

The allocation may use Poisson means below one; this is valid for
every $B>0$. The bound is uniform over the original rarity class.
Finite-level bounds can be optimized by taking an infimum over
finite constructions without asserting the existence of a single
infinite-level algorithm.

\section{One-unit null-work obstructions}\label{sec:lower-foundations}

This appendix proves the sign-testing and annular ingredients used in
Section~\ref{sec:matching-lower}. The weighted sign-testing mechanism is
an adaptation of Chen--Li--Qiao's gap-entropy argument, including their
Bernoulli extension \cite[Lemma 4.1 and Remark D.4]{clq}. We give the full
finite-gap version needed here, followed by the scale-sensitive gate and
its actual minimum-band-mass constraint.

\subsection{Weighted sign testing}

\begin{lemma}[Weighted sign entropy]\label{lem:weighted-sign}
Let $d_1>\cdots>d_K>0$, with $d_j/d_{j+1}\ge2$ and $d_1\le1/128$.
One observable binary sign test, with the same independent seed law at all
gaps, has error at most $1/10$ under every Bernoulli law $1/2\pm d_j$.
Suppose its stopping time $\tau$ has finite expectation under both signs at
every such gap.  Put $c_j=d_j^2\E_{-d_j}\tau$.
For every probability vector $\pi$ on these gaps,
\begin{equation}\label{eq:sc-weighted-sign}
 \sum_{j=1}^K\pi_jc_j\ge
 \frac{1+\operatorname{Ent}(\pi)}{257}.
\end{equation}
The stopping and decision rules are the same measurable functions of the
seed and observation tape at every gap; the realized gap is not supplied
as an input. The test may depend on the entire known finite gap list and
on the weights. Zero weights are allowed, with $0\log(1/0)=0$.
No termination or correctness assumption under the fair coin is required.
\end{lemma}

\begin{proof}
First, the ordinary stopped change-of-measure inequality gives a uniform
positive floor for each $c_j$.  The terminal sign decision has minus/plus
KL divergence at least
\[
 \operatorname{kl}(9/10,1/10)=\frac45\log9>\frac85.
\]
One Bernoulli observation has divergence at most
\[
 \frac{(2d_j)^2}{(1/2+d_j)(1/2-d_j)}\le\frac{64}{3}d_j^2.
\]
The stopped chain rule, justified by finite expected endpoint work, implies
$c_j>3/40>1/16$. More explicitly, apply the chain rule first to
the transcript capped at $N$. Encode a nonterminal decision by a third
symbol. Its probability tends to zero at both endpoints, so lower
semicontinuity of finite-alphabet KL and monotone convergence of
$\E_{-d_j}(\tau\wedge N)$ justify the stopped inequality.
This step uses no fair-law stopping assumption.

Set $b_0=1/1024$ and
\[
 n_j=\left\lfloor b_0/d_j^2\right\rfloor,\qquad
 u_j=\left\lceil4\E_{-d_j}\tau\right\rceil.
\]
Since $d_j\le1/128$, we have $n_j\ge b_0/(2d_j^2)\ge8$.
The KL divergence between the complete prefixes of length $n_j$ at the
minus and plus laws is at most $1/48$.  These prefixes include the
independent seed and unused bits if the test has already stopped.
Pinsker's inequality therefore gives total variation less than $1/8$.
An early minus is a plus-law error, so its probability under the minus law
is at most $1/10+1/8$.  Markov's inequality gives
$\Pp_{-d_j}(\tau>u_j)\le1/4$.  Consequently the event
\[
 E_j=\{\text{output minus},\ n_j<\tau\le u_j\}
\]
has $\Pp_{-d_j}(E_j)>1/4$.

Let $\Pp_0$ denote the fair-coin law with the same seed and policy.
Comparing deterministic prefixes of length $u_j$ gives KL at most
$4d_j^2u_j\le16c_j+1$.  Binary data processing and
$\operatorname{kl}(p,z)\ge p\log(1/z)-\log2$ yield
\begin{equation}\label{eq:sc-window-mass}
 \Pp_0(E_j)\ge\exp(-64c_j-8).
\end{equation}
Only deterministic finite prefixes have been compared with the fair law.

We next bound a weighted overlap pathwise.  On a path with finite terminal
time $t$, participation in $E_j$ implies
\[
 t>n_j\ge\frac{b_0}{2d_j^2},\qquad
 t\le4\E_{-d_j}\tau+1.
\]
In particular $t\ge2$ and
$c_j\ge(t-1)d_j^2/4\ge td_j^2/8$.
Order the participating gaps from smallest to largest.  Their squared
values grow by at least a factor of four, and the smallest has
$td_j^2>b_0/2$.  The corresponding costs are thus bounded below by
$b_0 4^r/16=4^{r-7}$, for $r=0,1,\ldots$.  Hence
\[
 \sum_j\frac{\1_{E_j}}{1+c_j}
 \le\sum_{r\ge0}\frac1{1+4^{r-7}}
 \le7+\sum_{r\ge7}4^{7-r}=\frac{25}{3}.
\]
On a nonterminating fair path no $E_j$ occurs, so the same inequality
holds.  Integrating and using \eqref{eq:sc-window-mass}, together with
$1+c_j\le e^{c_j}$, gives
\[
 Z:=\sum_je^{-65c_j}\le(25/3)e^8,
 \qquad \log Z<11.
\]
Gibbs' inequality for $\pi$ and the normalized masses $e^{-65c_j}/Z$
implies
\[
 \operatorname{Ent}(\pi)\le65\sum_j\pi_jc_j+\log Z.
\]
Writing $\bar c=\sum_j\pi_jc_j$ and using $\bar c\ge1/16$, we obtain
$1+\operatorname{Ent}(\pi)\le65\bar c+12\le257\bar c$, as required.
\end{proof}


\subsection{A moment comparison on the half-prior class}

Write $X_\rho=\{0\le q\le1/2:\int q\,d\mu\le\rho\}$.
For an integrable observable one-unit output $Y$ with finite uniform
variance and uniform bias at most $\rho/64$, put $V=\sup_{q\in X_\rho}\Var_qY$ and
$M=\sup_{q\in X_\rho}\E_qY^2$.

\begin{lemma}[Moment comparison]\label{lem:lf-moment}
For $0<\rho\le1/4$, one has $V\ge2\rho/3$ and $M\le3V/2$.
\end{lemma}
\begin{proof}
Let $f(d)=\E_{d,+}Y-\E_{d,-}Y$. Comparing the null population with
the constant prior $q=\rho$ gives $\int f\,d\mu\ge31/32$.
Conditional variance at this constant prior and Jensen's inequality yield
\[
 V\ge\rho(1-\rho)\int f^2\,d\mu
 \ge\rho(1-\rho)(31/32)^2\ge2\rho/3.
\]
Every mean has absolute value at most $65\rho/64$, so
$M\le V+(65\rho/64)^2\le3V/2$.
The constant prior also gives integrability of both fixed-sign second
moments; all conditional statements may be taken outside a common null set.
\end{proof}

\subsection{An activation gate that respects the magnitude scale}

The gate argument charges activation at the information
scale \(a^2\E_0T\).  In this subsection \(\mu\) is any probability law
supported on \((0,a]\), \(a\le1/4\), and
\[
 \mathcal X_\rho=
 \{q:0\le q\le\tfrac12,\ \textstyle\int q\,d\mu\le\rho\}.
\]
Let \((Y,T)\) be an observable one-coin procedure with uniform bias at
most \(\rho/64\), finite uniform variance \(V\), and finite expected work
under every \(q\in\mathcal X_\rho\).  Write
\[
 M_0=\sup_{q\in\mathcal X_\rho}\E_qY^2,\qquad
 W_0=\E_0T,\qquad P=M_0W_0.
\]
The moment comparison in Lemma~\ref{lem:lf-moment} gives
\(M_0\le3V/2\).

More explicitly, observability means that an independent seed $U$ and
the genuine observations generate $\mathcal F_n=\sigma(U,X_1,\ldots,X_n)$,
that $T$ is a nonnegative integer-valued stopping time, and that $Y$ is
$\mathcal F_T$-measurable. The procedure has no access to the latent
magnitude or sign. Restarts below refresh $U$ and the unused observation
block, while retaining that same latent state. These are the observation-model assumptions of
Section~\ref{sec:model}.


\begin{lemma}[A scale-sensitive baseline and gated activation]
\label{lem:information-gate}
Put
\[
 C_g=e^7,\qquad \kappa=64/3,\qquad
 c_g=(256\kappa C_g^4)^{-1}.
\]
Every procedure above satisfies
\begin{equation}
 P\ge\frac{\rho}{64a^2}.
 \label{eq:cross-baseline}
\end{equation}
If \(a^2P\le c_g\), there is an observable procedure \(Z\) on the same
one-coin experiment, with moment bound \(m_Z\), null work \(W_0^Z\), and
a measurable set \(\mathcal G\) of good magnitudes, such that
\begin{equation}
 \begin{gathered}
 \sup_{q\in\mathcal X_\rho}\E_qZ^2\le m_Z,\qquad
 \rho/16\le m_Z\le2C_g^4a^2P,\\
 m_ZW_0^Z\le C_g^4P,\qquad
 \mu(\mathcal G^c)\le144m_Z.
 \end{gathered}
 \label{eq:cross-gate-bounds}
\end{equation}
At every good gap, both signs have finite conditional expected work and
\begin{equation}
 \left|\E[Z\mid d,0]\right|\le\frac18,\qquad
 \E[Z\mid d,1]\ge\frac12,\qquad
 \E[Z^2\mid d,0]\le\frac{m_Z}{\rho},\qquad
 \E[Z^2\mid d,1]\le\frac{3m_Z}{\rho}.
 \label{eq:cross-gate-good}
\end{equation}
One common sign test with both errors below \(1/10\) on these gaps uses
\begin{equation}
 n=\left\lceil1024m_Z/\rho\right\rceil
 \le1040m_Z/\rho
 \label{eq:cross-gate-restarts}
\end{equation}
independent restarts of \(Z\) on unused observations of the same coin.
Its negative expected work is \(n\,\E[T_Z\mid d,0]\).
\end{lemma}

\begin{proof}
We first prove \eqref{eq:cross-baseline}.  Let \(P_-\) and \(P_+\) be
the full transcript laws with uniformly negative and uniformly positive
signs, respectively.  The priors \(q=0\) and \(q\equiv\rho\) give
\[
 \E_+Y-\E_-Y\ge31/32,\qquad
 \E_-Y^2\le M_0,\qquad \E_+Y^2\le M_0/\rho.
\]
The positive moment and expected work are finite because \(q\equiv\rho\)
is admissible.  Revealing \(D\) for the information calculation, and
using the stopped chain rule, gives
\[
 \mathfrak h(P_-,P_+)\le\tfrac12\kappa a^2W_0,
\]
where \(\mathfrak h=1-\operatorname{Aff}\).
Indeed the one-observation endpoint divergence is at most
\(\kappa d^2\) for \(d\le1/4\).
For any square-integrable output \(F\),
\begin{equation}
 |\E_PF-\E_QF|^2
 \le4(\E_PF^2+\E_QF^2)\mathfrak h(P,Q).
 \label{eq:cross-hellinger-moment}
\end{equation}
Thus
\[
 (31/32)^2
 \le(5\kappa/2)a^2P/\rho,
\]
which implies \eqref{eq:cross-baseline}.

For the gate, on a full observation tape define
\[
 S_n=\sum_{i=1}^n(2X_i-1),\qquad
 n_0=\lceil a^{-2}\rceil,\qquad H_0=\lceil a^{-1}\rceil,
\]
\begin{equation}
 G=\min\{n_0,\inf\{n\ge1:|S_n|\ge H_0\}\},\qquad
 \sigma=T\wedge G,\qquad E=\{T\ge G\}.
 \label{eq:cross-gate}
\end{equation}
The event \(E\) is measurable at \(\sigma\), including the tie \(T=G\).
Include the entire independent seed in the stopped prefix.
At \(\sigma\), \(|S_\sigma|\le H_0\) and \(\sigma\le n_0\).
For a fixed \(d\le a\) and sign \(s\), the likelihood relative to the
fair stopped-prefix law is
\[
 \exp\left\{(2s-1)S_\sigma\operatorname{arctanh}(2d)
       +\frac{\sigma}{2}\log(1-4d^2)\right\}.
\]
The inequalities
\[
 \operatorname{arctanh}(2d)\le8a/3,\qquad
 -\tfrac12\log(1-4d^2)\le8a^2/3
\]
show that the absolute logarithm is at most \(37/6<7\).
Every fixed-state prefix law, and the all-negative magnitude mixture,
therefore has density in \([C_g^{-1},C_g]\) relative to fair.

Let \(p=\Pp_{\rm fair}(E)\) and \(p_s(d)=\Pp_{d,s}(E)\).
The number $p$ is determined by the known policy, its seed law, and a
bounded fair-prefix experiment. It does not depend on the unknown prior
or genuine latent state. This is an existence argument; no efficient
computation of arbitrary policy probabilities is claimed. No termination
under the fair law beyond this bounded prefix is needed. Then
\begin{equation}
 p/C_g\le p_s(d)\le C_gp.
 \label{eq:cross-gate-probabilities}
\end{equation}
Bounded fair stopping gives
\(\E_{\rm fair}S_\sigma^2=\E_{\rm fair}\sigma\).
A gate activation either reaches \(n_0\) or hits \(H_0\), so
\begin{equation}
 p\le\left(n_0^{-1}+H_0^{-2}\right)\E_{\rm fair}\sigma
 \le2a^2\E_{\rm fair}\sigma
 \le2C_ga^2W_0.
 \label{eq:cross-gate-cost}
\end{equation}

Set \(Y_e=Y\1_{\{T<G\}}\), explicitly zero on \(E\).
This is a function of the bounded prefix.
Its second moment under either fixed state is at most \(C_g^2M_0\).
The stopped endpoint divergence and the same prefix domination give
\[
 \KL(P_{-d}^{\sigma},P_{+d}^{\sigma})
 \le\kappa d^2\E_{-d}\sigma
 \le\kappa C_g^2a^2W_0.
\]
Since \(\mathfrak h\le\KL/2\), \eqref{eq:cross-hellinger-moment} implies
\begin{equation}
 |\E_{+d}Y_e-\E_{-d}Y_e|^2
 \le4\kappa C_g^4a^2P.
 \label{eq:cross-early-output}
\end{equation}
Under \(a^2P\le c_g\), every such early mean difference is at most
\(1/8\) in absolute value.  In this branch \(p>0\): if \(p=0\), the
gate never activates under any state, so \(Y=Y_e\), contradicting
the integrated mean separation \(31/32\) from \(q=0,q\equiv\rho\).

Here is an actual observable procedure realizing the conditioning.
Keep the same genuine coin.  Repeatedly run the original procedure,
using a fresh original seed and unused iid flips for each trial.
Reject a trial if its original stopping time is before its gate.
When a trial reaches the gate, accept it, continue that original trial
to its original stopping time, and return
\[
 Z=C_gpY.
\]
Every rejected observation is charged.  Conditional on \(d,s\), trials
are iid with acceptance probability \(p_s(d)>0\).
Writing \(w_s(d)=\E[T\mid d,s]\), the renewal identity gives
\begin{equation}
 \E[T_Z\mid d,s]=\frac{w_s(d)}{p_s(d)}.
 \label{eq:cross-renewal}
\end{equation}
This identity includes the work of all rejected trials and the final
accepted trial. To verify it without assuming independence of acceptance
and trial cost, let $L$ be the first accepted trial. Conditional on $(d,s)$,
\[
 \E\sum_{k=1}^{L}T_k
 =\sum_{k\ge1}\E[T_k\mathbf1_{\{L\ge k\}}]
 =\sum_{k\ge1}(1-p_s(d))^{k-1}w_s(d)
 =w_s(d)/p_s(d).
\]
The event $\{L\ge k\}$ involves only preceding trials, so it is independent
of the current trial cost. It is not necessary that $T_k$ be independent
of its own acceptance event. Moreover,
\[
 \E[Z^2\mid d,s]
 =\frac{C_g^2p^2}{p_s(d)}\E[Y^2\1_E\mid d,s]
 \le C_g^3p\,\E[Y^2\mid d,s].
\]
With \(m_Z=C_g^3pM_0\), integration and
\eqref{eq:cross-gate-cost}--\eqref{eq:cross-renewal} give
\begin{equation}
 \sup_q\E_qZ^2\le m_Z,\qquad
 m_Z\le2C_g^4a^2P,\qquad
 W_0^Z\le C_gW_0/p,\qquad
 m_ZW_0^Z\le C_g^4P.
 \label{eq:cross-gate-moments}
\end{equation}
Expected work under every \(q\) is finite by the same comparison.

We establish conditional sign separation directly; conditioning has
not preserved the estimator's means.  Write
\[
 f_s(d)=\E[Y\mid d,s],\quad
 \ell_s(d)=\E[Y\1_E\mid d,s],\quad
 z_s(d)=\E[Z\mid d,s].
\]
The set \(\{f_+-f_-<3/4\}\) has mass below \(2\rho\).
Otherwise, writing that set as $S$ with $\mu(S)\ge2\rho$, the
prior $q=(\rho/\mu(S))\mathbf1_S$ lies in $\mathcal X_\rho$ and has total
mass $\rho$. It would contradict the bias guarantee and the mean
separation $31\rho/32$. This uses fractional prior values, so it applies
unchanged when $\mu$ has atoms.
Outside that set, \eqref{eq:cross-early-output} gives
\(\ell_+-\ell_-\ge5/8\).
Also
\[
 z_s(d)=\frac{C_gp}{p_s(d)}\ell_s(d),\qquad
 1\le\frac{C_gp}{p_s(d)}\le C_g^2.
\]
If \(|z_-(d)|\le1/8\), then \(|\ell_-(d)|\le1/8\),
\(\ell_+(d)\ge1/2\), and \(z_+(d)\ge1/2\).
The exceptional mass for \(|z_-|>1/8\) is at most \(64m_Z\).
Since
\[
 m_Z\le2C_g^4c_g=1/(128\kappa)<1/256,
\]
this preliminary mean-good set has mass at least
\(1-2\rho-64m_Z\ge1/4\).
Taking \(q\equiv\rho\) in the second-moment bound yields
\(m_Z\ge\rho/16\).

Finally, put \(a_Z(d)=\E[Z^2\mid d,0]\) and
\(b_Z(d)=\E[Z^2\mid d,1]\).
The set \(\{a_Z>m_Z/\rho\}\) has mass at most \(\rho\).
The set \(\{b_Z-a_Z>2m_Z/\rho\}\) has mass below \(2\rho\),
again by using a fractional admissible prior on a set of mass
at least \(2\rho\).  Removing these sets and null sets of infinite
conditional work gives total exceptional mass
\[
 5\rho+64m_Z\le144m_Z
\]
and proves \eqref{eq:cross-gate-good}.
Average \(n\) independent restarts of \(Z\) on unused blocks of the
same coin, and threshold at \(5/16\).
Both mean margins are at least \(3/16\), so Chebyshev gives errors
at most
\[
 \frac{3m_Z/\rho}{n(3/16)^2}\le1/12<1/10.
\]
The bound on \(n\) follows from \(m_Z\ge\rho/16\), and work adds,
proving \eqref{eq:cross-gate-restarts}.
\end{proof}

\subsection{A scaled annular obstruction}

For a positive vector \(v=(v_i)\), set
\[
 \mathcal G(v)=A[1+\Ent(v/A)],\qquad A=\sum_i v_i.
\]
A useful elementary identity is
\begin{equation}
 \frac{\partial\mathcal G}{\partial v_i}
 =1+\log(A/v_i)\ge1.
 \label{eq:cross-entropy-monotone}
\end{equation}
Thus this cost-entropy functional is coordinatewise increasing.
In particular, if \(u\le v_i\le U\) for \(K\) coordinates, then
\begin{equation}
 uK(1+\log K)\le\mathcal G(v)\le UK(1+\log K).
 \label{eq:cross-entropy-sandwich}
\end{equation}

\begin{proposition}[A gated annular lower bound]
\label{prop:scaled-annular}
Suppose \(\mu\) is supported on \((0,a]\), \(a\le1/4\).
Take a nonempty finite annular family
\[
 I_j=(b_j/2,b_j],\qquad b_j/b_{j+1}\ge4,\qquad
 b_1\le1/128,
\]
with positive masses \(w_j=\mu(I_j)\).  Define
\[
 w_*=\min_jw_j,\qquad
 A=\sum_jw_j/b_j^2,\qquad
 \pi_j=\frac{w_j/b_j^2}{A},\qquad
 \Gamma^{\circ}=A[1+\Ent(\pi)].
\]
Every one-coin procedure with the assumptions above satisfies
\begin{equation}
 V\E_0T\ge
 \frac{e^{-28}}{2\,000\,000}
       \min\{\rho\Gamma^{\circ},w_*/a^2\}.
 \label{eq:cross-scaled-annular}
\end{equation}
\end{proposition}

\begin{proof}
If \(a^2P\ge c_g\), use \(w_*\le1\) and \(M_0\le3V/2\).
Otherwise apply Lemma~\ref{lem:information-gate}.
If \(m_Z>w_*/288\), then
\[
 P>w_*/(576C_g^4a^2).
\]
If \(m_Z\le w_*/288\), at least half the mass of every annulus is
good.  Averaging selects \(d_j\in I_j\cap\mathcal G\) with
\[
 w_-^Z(d_j)\le\frac4{w_j}\int_{I_j}w_-^Z(d)\,\mu(dd).
\]
The selected gaps are separated by a factor greater than two and
bounded by \(1/128\). The representative points may depend on the
procedure being bounded. This is legitimate: the weighted sign lemma
holds for every such finite list. After the points have been selected,
the same already-constructed restart test is run at all of them; it is
not retuned to the selected point. The weights remain the normalized
band costs $\pi_j$, not the original band probabilities.
Apply Lemma~\ref{lem:weighted-sign} to the common restart test,
with weights \(\pi_j\):
\[
 \frac{1+\Ent(\pi)}{257}
 \le n\sum_j\pi_jd_j^2w_-^Z(d_j)
 \le\frac{4n}{A}W_0^Z.
\]
Consequently,
\[
 \rho\Gamma^{\circ}
 \le1028\rho nW_0^Z
 \le1\,069\,120\,m_ZW_0^Z
 \le1\,069\,120\,C_g^4P.
\]
Converting \(P=M_0W_0\) to \(VW_0\) completes all branches with the
displayed conservative constant.  No atomlessness is used.
\end{proof}

The entropy branch is available when $\rho a^2\Gamma^{\circ}\le w_*$.
This condition is invariant under a common rescaling of all magnitudes.

\section{Convex comparison, common pairs, and additive localization}
\label{sec:transfer-foundations}

This appendix proves the specialized transfer inputs used by the scalar
lower bound. The finite-experiment
convex-estimation principle has antecedents in
Juditsky--Nemirovski and Polyanskiy--Wu \cite{jn,pw};
the midpoint adaptive decomposition is due to Lee--Valiant \cite{lv}.
The arguments needed for this model are included in full.

\subsection{Finite categorical convex comparison}

\begin{lemma}[Finite categorical convex comparison]
\label{lem:finite-comparison}
Let $X$ be a nonempty compact convex subset of a finite-dimensional real vector
space. Let $\ell:X\to\R$ be affine, and let $p(x)$ be a strictly positive
categorical probability vector with affine continuous coordinates. For $n\ge1$
independent observations from $p(x)$, define
\[
 \Omega_n=\max\{\ell(x)-\ell(y):x,y\in X,
                  \operatorname{Aff}(p(x),p(y))^n\ge1/2\}.
\]
For every $\eta>0$, there is a bounded additive estimator
$F=c+\sum_{i=1}^n\phi_i(Z_i)$ satisfying
\begin{equation}
 \sup_{x\in X}\{\E_x(F-\ell(x))^2\}^{1/2}
 \le\tfrac32\Omega_n+\eta.
 \label{eq:finite-comparison}
\end{equation}
\end{lemma}

\begin{proof}
We reproduce the short argument needed for the numerical constant. Set
$s=\log2$, let $\mathcal F$ be the finite-dimensional space of additive
functions, and define, for $\alpha>0$,
\[
 \Phi(x,y;\phi,\alpha)=\ell(x)-\ell(y)
 +\alpha\log\E_y e^{\phi/\alpha}
 +\alpha\log\E_x e^{-\phi/\alpha}+2\alpha s.
\]
The function is continuous and concave in $(x,y)$: independence writes each
log moment as a sum of logarithms of positive affine forms. It is convex in
$(\phi,\alpha)$ by the perspective transformation of a log moment generating
function. Sion's minimax theorem \cite[corollary following Theorem 3.4, p.~174]{sion} therefore applies to $X\times X$ and
$\mathcal F\times(0,\infty)$. Only $X\times X$ needs to be compact;
the open, unbounded $\alpha$ domain is convex and causes no problem in
Sion's one-compact-set form. The infimum in $(\phi,\alpha)$ need not be
attained, and the later argument uses a strict slack $R>S$. Write
$2S=\inf_{\phi,\alpha>0}\max_{x,y}\Phi$.

For a fixed pair, the half log likelihood ratio belongs to $\mathcal F$ and
Cauchy--Schwarz gives
\[
 \inf_{\psi\in\mathcal F}
 \bigl[\log\E_y e^\psi+\log\E_x e^{-\psi}\bigr]
 =2\log\operatorname{Aff}(p(x)^{\otimes n},p(y)^{\otimes n}).
\]
Substituting $\phi=\alpha\psi$ and minimizing over $\alpha>0$ gives the value
$\ell(x)-\ell(y)$ when the affinity is at least $1/2$, and $-\infty$ otherwise.
Consequently $2S=\Omega_n$. The constrained maximum exists because its feasible
set is nonempty and compact.

For $R>S$, choose $(\phi,\alpha)$ with $\max\Phi\le2R$. On the diagonal,
Cauchy--Schwarz implies $\Phi(x,x;\phi,\alpha)\ge2\alpha s$, hence
$\alpha\le R/s$. Set
\[
 U=\max_x\{\ell(x)+\alpha\log\E_xe^{-\phi/\alpha}\},\qquad
 V=\max_x\{-\ell(x)+\alpha\log\E_xe^{\phi/\alpha}\}.
\]
Then $U+V+2\alpha s\le2R$. The estimator
$F=\phi+(U-V)/2$ satisfies
\[
 \E_xe^{\pm(F-\ell(x))/\alpha}\le e^{R/\alpha-s}.
\]
Therefore $\Pp_x\{|F-\ell(x)|>R+t\}\le e^{-t/\alpha}$ for $t\ge0$.
Tail integration yields
\[
 \E_x(F-\ell(x))^2
 \le R^2+2R\alpha+2\alpha^2
 \le R^2\left(1+\frac2{\log2}+\frac2{(\log2)^2}\right)<9R^2.
\]
Let $R\downarrow S$, with the slack chosen for $\eta$. Every component is
bounded because its alphabet is finite. This proves
\eqref{eq:finite-comparison}.
\end{proof}


\subsection{Midpoint decomposition for interleaved observations}

We use $\mathfrak h(P,Q)=1-\Aff(P,Q)$, so $0\le\mathfrak h\le1$ and
$\operatorname{TV}(P,Q)\le\sqrt{2\mathfrak h(P,Q)}$.
The following finite-prefix argument is the midpoint decomposition of
Lee--Valiant \cite[Lemma 17]{lv}, written in the form needed here.

\begin{lemma}[Midpoint work decomposition]\label{lem:tf-midpoint}
Fix two endpoint priors $q_-,q_+$ and their midpoint
$\bar q=(q_-+q_+)/2$. Suppose every capped randomized one-unit policy
$A$ obeys
\[
 \mathfrak h(P_-^A,P_+^A)\le\kappa\,\E_{\bar q}T_A.
\]
Then every population algorithm with a deterministic total cap satisfies
\[
 \mathfrak h(P_-^{\rm pop},P_+^{\rm pop})
 \le\kappa\,\E_{\bar q}T_{\rm pop}.
\]
It suffices that the local hypothesis hold up to that same total cap.
If instead $q_\pm\le1/2$ and the local bound uses null work
$\lambda\E_0T_A$, the conclusion holds with $\kappa=2\lambda$.
The statement also holds for finitely many observable component types,
each with its own endpoint pair and midpoint, when the same coefficient
bounds every local policy for every type. Types may be selected adaptively, but each fresh unit's type is fixed
before its first observation and retained on every revisit.
\end{lemma}

\begin{proof}
Fix total cap $K$ and condition on the algorithm's independent seed.
Defer latent draws until a unit is first observed, which changes no
observable law. Index units by that first observation. A consistent global leaf $z$
contains local binary words $h_i(z)$, with empty words for unused slots.
The endpoint leaf densities factor as
\[
 p_\pm^{\rm pop}(z)=\prod_{i=1}^K p_\pm(h_i(z)).
\]
The common policy-consistency indicator is implicit. Put
$m_i=(p_-(h_i)+p_+(h_i))/2=p_{\bar q}(h_i)$ and
$c_i=\sqrt{p_-(h_i)p_+(h_i)}$. Then
\[
 \prod_i m_i-\prod_i c_i
 \le\sum_i(m_i-c_i)\prod_{j\ne i}m_j.
\]
The integral of $\prod_i m_i$ over consistent leaves is one, since it
is the population law in which every fresh unit has midpoint prior.
Integrating the inequality therefore bounds population squared Hellinger
by a sum of hybrid squared Hellinger distances. In hybrid $i$ only
unit $i$ has an endpoint prior; all other units have midpoint priors.

Simulate all other units from independent midpoint tapes included in
an auxiliary seed, and run the population rule with one genuine unit
in slot $i$. This is a capped local policy. Its full transcript,
including that seed, determines the hybrid global transcript.
Data processing bounds hybrid distance by local distance. The expected
midpoint local works sum to $\E_{\bar q}T_{\rm pop}$ because each
purchased observation belongs to exactly one slot. This proves the
first assertion, including after integrating the original seed.

If $q_\pm\le1/2$, the midpoint latent law dominates half the null
latent law. The same measure inequality holds after passing through
any observation policy and auxiliary seed. Thus
$\E_0T_A\le2\E_{\bar q}T_A$.

For observable types, the type selected for a slot is a function of the
history before its first observation. Use the resulting type-specific
kernel in each leaf product. To implement hybrid slot $i$, first generate
all observations before that slot under the componentwise midpoint law.
Conditional on this prehistory and its selected type, the genuine unit
has that type's prescribed endpoint law and is fresh, independently of
the remaining auxiliary randomness. This is conditional independence:
the latent law may depend on the type contained in the seed.

Retain the already generated tapes of other units on every revisit.
For later slots, an independent array of potential tapes indexed by slot
and type can be placed in the auxiliary seed; choose the appropriate
unused tape only after that slot's type has been selected. The seed's
distribution is the same under the two endpoint laws for the genuine
unit. Conditional on the prehistory, the local hypothesis for its selected
type applies. If slot $i$ is never reached, the genuine local work is zero.
Summing and integrating counts every purchased observation once under
the actual componentwise midpoint population. No latent magnitude or
sign of the genuine unit is revealed to the local policy.
\end{proof}

For unbounded stopping in the applications here, the same endpoint pair
works at every cap, and the original uniform budget gives expected work
at most $B$ at both endpoints and at their admissible midpoint. Let $N$
be original purchased work, and cap after $L$ observations, recording an
overflow flag if needed. Any terminal event restricted to $\{N\le L\}$
is measurable from this capped transcript. Therefore
\[
 \operatorname{TV}(P_-,P_+)
 \le\operatorname{TV}(P_-^{(L)},P_+^{(L)})
       +\Pp_-(N>L)+\Pp_+(N>L)
 \le\sqrt{2\kappa B}+2B/L.
\]
Letting $L\to\infty$ gives the terminal testing bound. For an oracle
simulation, cap at the original observation count; its genuine work
is no larger, and the same argument applies. This avoids exchanging a
cap limit with a separately chosen pair.

\subsection{The half-prior class and null work}

Put
\[
 X_\rho=\left\{q:0\le q\le\tfrac12,\quad
                  \theta(q):=\int q\,d\mu\le\rho\right\},
 \qquad 0<\rho\le\tfrac14.
\]
For a capped observable one-coin procedure with bounded real output \(Y\),
write
\[
 V=\sup_{q\in X_\rho}\Var_qY,\qquad
 M=\sup_{q\in X_\rho}\E_qY^2,\qquad W_0=\E_0T,
\]
and define
\begin{equation}
 \mathfrak J_0(b)=\inf\left\{VW_0:
       \sup_{q\in X_\rho}|\E_qY-\theta(q)|\le b,
       \ T\text{ capped},\ Y\text{ bounded}\right\}.
 \label{eq:null-bias-complexity}
\end{equation}
An empty infimum is infinity.  The function is nonincreasing in \(b\),
and \(\mathfrak J_0(\rho/2)=0\).

\subsection{Transfer at an arbitrary bias tolerance}

\begin{theorem}[Bias-indexed common pairs and population lower bounds]
\label{thm:bias-transfer}
Suppose that every capped bounded one-coin output with uniform bias at
most \(b>0\) over \(X_\rho\) satisfies \(VW_0\ge J\), where
\(0<J<\infty\).  Then, for every \(0<\delta\le b/3\), there is one pair
\(q_-,q_+\in X_\rho\) with target separation at least \(\delta\) such
that simultaneously for all capped randomized local policies \(A\),
\begin{equation}
 \mathfrak h(P_{q_-}^A,P_{q_+}^A)
 \le42\frac{\delta^2}{J}\E_0T_A.
 \label{eq:bias-common-pair}
\end{equation}
Consequently,
\begin{equation}
 R_{\mu,\rho}(B)\ge\frac1{128}
                 \min\{b,\sqrt{J/B}\}.
 \label{eq:bias-population-transfer}
\end{equation}
The conclusion includes fully adaptive creation, revisiting, interleaving,
nonlinear outputs and unbounded stopping.
\end{theorem}

\begin{proof}
The premise forces \(b<\rho/2\), since the zero-work constant \(\rho/2\)
would otherwise contradict it.  First fix a finite strictly positive
categorical statistic of a capped policy of null work \(w>0\).
For \(n\) independent runs, let
\(F=c+\sum_{i=1}^nY_i\) be a bounded additive estimator of worst RMS
\(r\).  Put \(b_i=\E_0Y_i\), \(b_F=\E_0F\), and choose an independent
uniform index \(I\).  Running just its component and returning
\(Z=n(Y_I-b_I)+b_F\) preserves all population means at null work \(w\).

The laws under \(q\) and zero share a subprobability of mass
\(1-\theta(q)\), obtained by weighting the negative-sign kernel by
\(1-q(d)\).  The elementary common-mixture variance inequality gives
\[
 (\E_qY_i-b_i)^2\le
 \frac{2\theta(q)}{1-\theta(q)}
                 (\Var_qY_i+\Var_0Y_i).
\]
Indeed, either mean's displacement from the normalized common part is
at most the square root of \(\theta/(1-\theta)\) times its standard
deviation; square the sum of the two displacements.  Independence in
\(F\) now yields
\[
 \sup_q\Var_qZ\le
 nr^2\left(1+\frac{4\rho}{1-\rho}\right)\le\tfrac73nr^2.
\]
If \(r<b\), apply the assumed obstruction to \(Z\); if \(r\ge b\),
retain that inequality.  In both cases,
\[
 r\ge\min\{b,\sqrt{3J/(7nw)}\}.
\]

Lemma~\ref{lem:finite-comparison} gives an additive estimator of
RMS at most \((3/2)\Omega_n+\eta\), where \(\Omega_n\) is the maximum
target separation whose product affinity is at least \(1/2\).
Letting \(\eta\downarrow0\) therefore gives
\[
 \Omega_n\ge\tfrac23\min\{b,\sqrt{3J/(7nw)}\}.
\]
Set \(x=J/(21\delta^2w)\).  If \(x\ge2\), take
\(n=\lfloor x\rfloor\); then \(\Omega_n\ge2\delta\), and the
corresponding pair satisfies
\[
 \mathfrak h\le\frac{\log2}{n}
 \le42\log2\,\frac{\delta^2w}{J}
 \le42\frac{\delta^2w}{J}.
\]
If \(x<2\), the last right side exceeds one, so any separated pair
suffices.  If \(w=0\), the seed-only first-observation event has
probability zero; all local laws agree and \(q_-=0,q_+=\delta\)
suffices.

We now obtain one pair simultaneously for every capped local policy.
For an arbitrary capped policy, partition its full transcript into finitely
many measurable cells, then replace the cell independently by a uniform
cell with probability \(0<\varepsilon<1\).  Run the policy even on
replacement outcomes, so its null work is unchanged.  The resulting
strictly positive category probabilities and the target have compact
convex finite-dimensional range, so the preceding comparison applies.
For any finite collection of these experiments, independently mix them
and reveal the chosen index. Categories with zero mixture weight can be
dropped, leaving a strictly positive finite alphabet. For a fixed one-unit
policy, its categorical law is affine in $q$: condition on the one genuine
latent state first. This affineness is not being asserted for an entire
adaptive population transcript.

Set $k=42\delta^2/J$ and use the oriented compact convex pair set
\[
 K_\delta=\{(q_-,q_+)\in X_\rho^2:
                 \theta(q_+)-\theta(q_-)\ge\delta\}.
\]
For experiment $i$, let $w_i$ be its null work and
$f_i(q_-,q_+)=\mathfrak h(P_{q_-}^{A_i},P_{q_+}^{A_i})-kw_i$.
Each $f_i$ is continuous and convex on $K_\delta$.
For any mixture weights $\lambda$, the revealed index gives exactly
$\mathfrak h_{\rm mix}=\sum_i\lambda_i\mathfrak h_i$ and $w_{\rm mix}=\sum_i\lambda_iw_i$.
The preceding one-experiment argument therefore implies
\[
 \min_{z\in K_\delta}\sum_i\lambda_i f_i(z)\le0
 \qquad\text{for every }\lambda\text{ in the finite simplex}.
\]
Applying convex minimax in its indicated direction gives
\[
 \min_{z\in K_\delta}\max_i f_i(z)
 =\max_\lambda\min_{z\in K_\delta}\sum_i\lambda_i f_i(z)\le0.
\]
The left minimum is attained. Thus the finite family has a single feasible
pair; separate pairs for the individual experiments would not suffice.
All its constraint sets are closed in $K_\delta$.
The finite intersection property supplies one pair for all experiments.
The intersection ranges over all caps, finite partitions and positive
noise levels at once, so this pair does not depend on the stopping cap.
The pair is fixed while $\varepsilon\downarrow0$; finite categorical
Hellinger distance is continuous even when some limiting cells have zero
mass. Full squared Hellinger distance is the supremum over finite
measurable partitions. For completeness, put $R=P+Q$ and $t=dP/dR$.
Partition the range of $t$ into intervals of width at most $1/m$.
Since $g(t)=\sqrt{t(1-t)}$ satisfies $|g(u)-g(v)|\le\sqrt{|u-v|}$,
the full and coarsened affinities differ by at most $2/\sqrt m$.
Thus the supremum statement holds without a countable-generation
assumption on the transcript space. This proves
\eqref{eq:bias-common-pair} for arbitrary transcript spaces.
The quantifier is one pair for every policy and cap at fixed
$(\mu,\rho,b,J,\delta)$. The pair may change when $\delta$, and hence
the chosen budget-dependent separation, changes.

The midpoint decomposition in Lemma~\ref{lem:tf-midpoint}
bounds population squared Hellinger by
\(84\delta^2B/J\): the local null cost is at most twice the midpoint
cost, and summing midpoint costs counts each observation once.  Set
\[
 \delta=\frac1{32}\min\{b,\sqrt{J/B}\}.
\]
The distance is at most \(84/1024<1/8\), hence TV is at most \(1/2\).
The same pair works for every deterministic observation cap.  Finite
expected endpoint work and increasing capped prefixes give the same
terminal TV bound for unbounded stopping.  Thresholding an estimator at
the target midpoint gives RMS at least \(\delta/4\), proving
\eqref{eq:bias-population-transfer}.  A budget over the full rarity class
implies the required budget over \(X_\rho\).
\end{proof}


\subsection{A direct sum of conditional experiments}

Let $\sigma_1,\ldots,\sigma_m$ be nonzero finite submeasures with
$\sum_i\sigma_i\le\mu$. Write $\tau_i=\sigma_i((0,1/2])$ and
$\nu_i=\sigma_i/\tau_i$. Disjoint restrictions are a special case;
fractional restrictions of atoms are also allowed. Choose
\[
 0<r_i\le1/4,\qquad \sum_i\tau_i r_i\le\rho.
\]
Within component $i$, use the half-prior subclass
\[
 X_{r_i}(\nu_i)=\{0\le q\le1/2:\ \textstyle\int q\,d\nu_i\le r_i\}.
\]
The original population lower bound is valid on this subclass of the full
rarity class.

\begin{theorem}[Additive localization]\label{thm:additive-localization}
Suppose that, for each $i$, every capped bounded observable one-unit output
with uniform bias at most $b_i>0$ over $X_{r_i}(\nu_i)$ satisfies
\[
 \left(\sup_q\Var_qY\right)\E_{i,0}T\ge J_i>0.
\]
Then, for every $B>0$,
\begin{equation}\label{eq:additive-localization}
 \boxed{\quad
 R_{\mu,\rho}(B)\ge\frac1{128}
 \sum_{i=1}^m\tau_i\min\{b_i,\sqrt{J_i/B}\}.
 \quad}
\end{equation}
The same total budget $B$ occurs in every summand.
\end{theorem}

\begin{proof}
Theorem~\ref{thm:bias-transfer} applies within every component. Set
\[
 \delta_i=\frac1{32}\min\{b_i,\sqrt{J_i/B}\}.
\]
It supplies a pair $q_{i,-},q_{i,+}\in X_{r_i}(\nu_i)$, oriented to have
conditional target separation at least $\delta_i$, such that simultaneously
for every capped local policy $A$,
\begin{equation}\label{eq:additive-local-pair}
 \mathfrak h(P_{i,-}^A,P_{i,+}^A)
 \le42\frac{\delta_i^2}{J_i}\E_{i,0}T_A
 \le\frac{42}{1024B}\E_{i,0}T_A.
\end{equation}

Fix an arbitrary original algorithm with uniform expected budget $B$.
Consider an oracle experiment that can select a component at no cost and
then receive a genuine coin with magnitude law $\nu_i$. Its component label
is visible, while its magnitude and sign remain hidden. The algorithm can
adaptively select, revisit, and interleave these coins. Apply the argument
to the simulation of the fixed original algorithm described below. At the
componentwise midpoint it induces the original parameter
$\bar q=(q_-+q_+)/2\in\mathcal Q_\rho$, so its genuine midpoint work is
at most the original work and hence at most $B$. Under the conditional
midpoint prior, $q\le1/2$, so the midpoint law dominates half the null law.
Lemma~\ref{lem:tf-midpoint} therefore bounds terminal squared
Hellinger by $84/1024<1/8$: the local coefficient in
\eqref{eq:additive-local-pair} is common to all components, and the summed
conditional midpoint costs count each purchased observation once.
Increasing capped prefixes handles unbounded stopping. The total target
separation is at least $\sum_i\tau_i\delta_i$.

Every original algorithm can be simulated in this stronger experiment.
Privately choose the component label with probabilities $\tau_i$ whenever
it creates a fresh coin. On the remaining probability mass, simulate the
known all-negative residual law for free. Keep the chosen label when the
coin is revisited. For fractional submeasures, integrating this private
mixture label gives the original admissible parameters
\[
 q_\pm(d)=\sum_i\frac{d\sigma_i}{d\mu}(d)q_{i,\pm}(d),
 \qquad 0\le q_\pm\le1/2,\qquad
 \int q_\pm\,d\mu\le\sum_i\tau_i r_i\le\rho.
\]
The number of genuine oracle observations does not exceed the original
number. Thus the oracle lower bound applies to the original experiment.
The Hellinger bound gives TV at most $1/2$, and midpoint testing gives RMS
at least one quarter of the target separation. This is
\eqref{eq:additive-localization}.
\end{proof}

For a direct explicit input, suppose component $i$ is supported in
$(0,a_i]$, $a_i\le1/4$, and select a finite separated family of small-gap
annuli. Let $\Gamma_i^{\circ}$ and $w_{*,i}$ denote its conditional cost entropy
and minimum band mass. Proposition~\ref{prop:scaled-annular} gives
\[
 b_i=r_i/64,\qquad
 J_i=c_{\mathrm{ann}}
       \min\{r_i\Gamma_i^{\circ},w_{*,i}/a_i^2\},\qquad
 c_{\mathrm{ann}}=\frac{e^{-28}}{2{,}000{,}000}.
\]
Consequently, writing the corresponding unnormalized quantities with
respect to $\sigma_i$ as $\Gamma_{\sigma_i}^{\circ}=\tau_i\Gamma_i^{\circ}$ and
$w_{*,\sigma_i}=\tau_iw_{*,i}$, we obtain
\begin{equation}\label{eq:additive-annular}
 R_{\mu,\rho}(B)\gtrsim
 \sum_i\min\left\{\tau_i r_i,
 \sqrt{\frac{
 \min\{\tau_i r_i\Gamma_{\sigma_i}^{\circ},
           \tau_iw_{*,\sigma_i}/a_i^2\}}B}\right\}.
\end{equation}
For clarity, the squared component mass is essential before cancellation:
\[
 \tau_i^2J_i=c_{\mathrm{ann}}
 \min\{\tau_i r_i\Gamma_{\sigma_i}^{\circ},
             \tau_iw_{*,\sigma_i}/a_i^2\}.
\]
This explains the extra $\tau_i$ in the saturation term of
\eqref{eq:additive-annular}; it must not be dropped when passing between
conditional and unnormalized quantities.
Section~\ref{sec:matching-lower} applies this bound to the components
constructed by the quantile-core filter and packing.

\clearpage
\phantomsection
\addcontentsline{toc}{section}{References}
\begin{thebibliography}{9}
\bibitem{lv}
J. C. H. Lee and P. Valiant.
\emph{Uncertainty about Uncertainty: Optimal Adaptive Algorithms for
Estimating Mixtures of Unknown Coins}.
Proceedings of the 2021 ACM--SIAM Symposium on Discrete Algorithms
(SODA), 414--433, 2021.
\url{https://arxiv.org/abs/1904.09228v3}.

\bibitem{clq}
L. Chen, J. Li, and M. Qiao.
\emph{Towards Instance Optimal Bounds for Best Arm Identification}.
Proceedings of the 2017 Conference on Learning Theory,
Proceedings of Machine Learning Research 65:535--592, 2017.
\url{https://proceedings.mlr.press/v65/chen17b.html}.

\bibitem{jn}
A. B. Juditsky and A. S. Nemirovski.
\emph{Nonparametric estimation by convex programming}.
Annals of Statistics 37(5A):2278--2300, 2009.
\url{https://doi.org/10.1214/08-AOS654}.

\bibitem{pw}
Y. Polyanskiy and Y. Wu.
\emph{Dualizing Le Cam's method for functional estimation, with
applications to estimating the unseens}.
arXiv:1902.05616v2, 2021.
\url{https://arxiv.org/abs/1902.05616v2}.

\bibitem{hoeffding}
W. Hoeffding.
\emph{Probability Inequalities for Sums of Bounded Random Variables}.
Journal of the American Statistical Association 58(301):13--30, 1963.
\url{https://doi.org/10.1080/01621459.1963.10500830}.

\bibitem{sion}
M. Sion.
\emph{On general minimax theorems}.
Pacific Journal of Mathematics 8(1):171--176, 1958.
\url{https://doi.org/10.2140/pjm.1958.8.171}.
\end{thebibliography}

\end{document}
